\exercisesection{}

在下面的习题中，当 E 是巴拿赫空间时，我们用 \(\mathcal{C}alk(E)\) 表示卡尔金代数 \(\mathcal{L}(E)/\mathcal{L}^{c}(E)\) 。用 \(\pi\) 表示从 \(\mathcal{L}(E)\) 到 \(\mathcal{C}alk(E)\) 上的典范投影 。

若 E 是复希尔伯特空间，我们给 \(\mathcal{C}alk(\mathrm{E})\) 赋以由 \(\mathcal{L}(\mathrm{E})\) 的对合经取商而诱导的对合。

\begin{enumerate}
\def\labelenumi{\arabic{enumi})}
\tightlist
\item
  设 E 是局部凸空间。
\end{enumerate}

\begin{enumerate}
\def\labelenumi{\alph{enumi})}
\item
  假设 E 中 0 的每个邻域都包含 E 的一个有限余维数的向量子空间。证明：E 的弗雷德霍姆自同态所成的空间 \(\mathcal{F}(\mathrm{E})\) 在赋以有界收敛拓扑的 \(\mathcal{L}(\mathrm{E})\) 中不是闭的。
\item
  给出一个无穷维弗雷歇空间的例子，使得其中 0 的每个邻域都包含 E 的一个有限余维数的向量子空间。c) 设 \(E_{\sigma}\) 是赋以弱拓扑的无穷维巴拿赫空间。集合 \(\mathcal{F}(E_{\sigma})\) 在 \(\mathcal{L}(E_{\sigma})\) 中不是开的。
\end{enumerate}

\begin{enumerate}
\def\labelenumi{\arabic{enumi})}
\setcounter{enumi}{1}
\tightlist
\item
  设 E 是无穷维巴拿赫空间，且 \(u \in \mathcal{L}(\mathrm{E})\) 。我们用 \(\mathrm{S}(u)\) 表示如下的复数 \(\lambda\) 所成的集：存在实数 \(\delta > 0\) 和 E 的一个有限余维数子空间 F，使得 \(\| u(x) - \lambda x \| \geqslant \delta \| x \|\) 对所有 \(x \in \mathrm{F}\) 成立。
\end{enumerate}

\begin{enumerate}
\def\labelenumi{\alph{enumi})}
\item
  集合 \(S(u)\) 在 C 中是开的，且包含 u 的谱在 C 中的补集。
\item
  若 \(\lambda \in \mathrm{S}(u)\) ，则自同态 \(u - \lambda 1_{\mathrm{E}}\) 具有闭的像和有限维的核。
\item
  补集 \(\mathrm{F}(u) = \mathbf{C} - \mathrm{S}(u)\) 非空。（证明 \(\mathrm{Sp}_{\mathrm{e}}(u)\) 中一个模最大的元素属于 \(\mathrm{F}(u)\) 。）
\end{enumerate}

\begin{enumerate}
\def\labelenumi{\arabic{enumi})}
\setcounter{enumi}{2}
\tightlist
\item
  设 E 是无穷维复希尔伯特空间。
\end{enumerate}

\begin{enumerate}
\def\labelenumi{\alph{enumi})}
\item
  对合代数 \(\mathcal{C}alk(E)\) 是星代数。
\item
  设 u 是 \(\mathcal{L}(\mathrm{E})\) 的使 \(u^{*}u\) 为紧算子的元素。证明 u 是紧算子。
\item
  设 G 是 \(\mathcal{C}alk(\mathrm{E})\) 中可逆元素所成的群，\(G_{0}\) 是它的单位分量。设 \(p: G \to G/G_{0}\) 是典范投影。存在唯一的群同构 \(\alpha: Z \to G/G_{0}\) ，使得对 E 的每个弗雷德霍姆自同态 u 都有 \(\alpha(\operatorname{ind}(u)) = p(\pi(u))\) 。
\end{enumerate}

\begin{enumerate}
\def\labelenumi{\arabic{enumi})}
\setcounter{enumi}{3}
\tightlist
\item
  设 E 是复希尔伯特空间。设 \(\xi\) 是 \(\mathcal{C}alk(E)\) 的一个可逆元素，我们称 \(\xi\) 的指标为 E 的任一弗雷德霍姆自同态 \(u\) （满足 \(\pi(u) = \xi\) ）的指标。
\end{enumerate}

\begin{enumerate}
\def\labelenumi{\alph{enumi})}
\item
  证明：对每个埃尔米特的 \(\xi \in \mathcal{C}alk(\mathrm{E})\) ，存在埃尔米特的 \(u \in \mathcal{L}(\mathrm{E})\) 使得 \(\pi(u) = \xi\) 。
\item
  证明：对每个酉的且指标为零的 \(\xi \in \mathcal{C}alk(\mathrm{E})\) ，存在酉的 \(u \in \mathcal{L}(\mathrm{E})\) 使得 \(\pi(u) = \xi\) 。（设 \(v\) 可逆且满足 \(\pi(v) = \xi\) ；考察 \(v = j|v|\) 这一 \(v\) 的极分解，验证 \(|v| - 1_{\mathrm{E}}\) 是紧的，再推出 \(v - j\) 是紧的。）
\item
  证明：对每个酉的 \(\xi \in \mathcal{C}alk(\mathrm{E})\) （满足 \(\xi^2 = 1\) ），存在酉的 \(u \in \mathcal{L}(\mathrm{E})\) 使得 \(u^2 = 1\) 且 \(\pi(u) = \xi\) 。
\end{enumerate}

\(d)\) 对每个元素 \(p \in \mathcal{C}alk(\mathrm{E})\) ，只要 \(p^2 = p\) ，就存在投影 \(\mathrm{P} \in \mathcal{L}(\mathrm{E})\) 使得 \(\pi(\mathrm{P}) = p\) 。

\begin{enumerate}
\def\labelenumi{\arabic{enumi})}
\setcounter{enumi}{4}
\tightlist
\item
  设 E 是复希尔伯特空间。
\end{enumerate}

\begin{enumerate}
\def\labelenumi{\alph{enumi})}
\tightlist
\item
  设 \(n\) 是满足 \(n \geqslant 2\) 的整数。我们置 \(\mathrm{I} = \mathbf{N} \times \{0, \ldots, n\}\) 。设 \((\sigma_i)_{1 \leqslant i \leqslant n}\) 是 I 的由下式定义的置换族
\end{enumerate}

\[
\begin{array}{c} \sigma_ {i} (m, i) = (m - 1, i) \text {for} m \geqslant 1 \\ \sigma_ {i} (0, i) = (0, i - 1) \\ \sigma_ {i} (m, i - 1) = (m + 1, i - 1) \text {for} m \geqslant 0 \\ \sigma_ {i} (m, j) = (m, j) \text {for} m \geqslant 0 \text {and} j \notin \{i, i - 1 \}. \end{array}
\]

\begin{enumerate}
\def\labelenumi{\alph{enumi})}
\setcounter{enumi}{1}
\item
  用 \((e_{m,i})\) 表示 \(\ell^{2}(\mathrm{I})\) 的典范标准正交基。设 \(u_{i}\) 是 \(\ell^{2}(\mathrm{I})\) 的自同态，使得 \(u_{i}(e_{m,j}) = e_{\sigma_{i}(m,j)}\) 对所有 \((m,j) \in \mathrm{I}\) 成立。自同态 \(u_{i}\) 是酉的。
\item
  元素 \(\pi(u_i)\) 两两交换。
\end{enumerate}

\(d)\) 不存在 E 的由两两交换的酉自同态组成的族 \((v_{1},\ldots,v_{n})\) ，使得 \(\pi(v_{i})=\pi(u_{i})\) 对所有 \(i\) 成立。

\begin{enumerate}
\def\labelenumi{\arabic{enumi})}
\setcounter{enumi}{5}
\tightlist
\item
  设 E 是复希尔伯特空间。回顾 \(\mathcal{C}alk(E)\) 是星代数（习题 3）。
\end{enumerate}

  设 \(\mathfrak{F}\) 是 \(\mathbf{N}\) 上的一个非主超滤子。设 F 是 E 中沿 \(\mathfrak{F}\) 弱收敛于 0 的有界序列所成的向量空间，\(F_0\) 是 F 中由沿 \(\mathfrak{F}\) 依范数收敛于 0 的序列组成的子空间。我们置 \(\widetilde{\mathrm{H}} = \mathrm{F} / \mathrm{F}_0\) 。a) 对 F 的所有元素 \(x = (x_n)\) 和 \(y = (y_n)\) ，证明序列 \((\langle x_n | y_n \rangle)_{n \in \mathbf{N}}\) 沿 \(\mathfrak{F}\) 收敛。用 \(\langle x | y \rangle_{\mathrm{F}}\) 表示其极限。

\begin{enumerate}
\def\labelenumi{\alph{enumi})}
\setcounter{enumi}{1}
\item
  证明 \((x,y)\mapsto\langle x\mid y\rangle_{\mathrm{F}}\) 是 F 上的正埃尔米特形式，且 \(\langle x\mid x\rangle_{F}=0\) 当且仅当 \(x\in F_{0}\) ；因而这个埃尔米特形式在 \(\widetilde{H}\) 上定义了一个分离的预希尔伯特空间结构。
\item
  设 H 是由分离的预希尔伯特空间 \(\widetilde{H}\) 完备化得到的希尔伯特空间。证明：对 \(\mathcal{L}(E)\) 中的每个 u，都得到一个元素 \(\varrho(u)\) （属于 \(\mathcal{L}(H)\) ），这只需把 F 到 F 的线性映射 \((x_{n})\mapsto(u(x_{n}))\) 过渡到商。d) 映射 \(u\mapsto\varrho(u)\) 诱导 C*-代数的一个等距态射 \(\mathcal{C}alk(E)\to\mathcal{L}(H)\) 。
\end{enumerate}

（我们以后将看到，每个 C*-代数 A 都有一个等距表示 \(A \to \mathcal{L}(E)\) ，其中 E 是某个希尔伯特空间，参见 V, p. 442 的定理 2。）

\begin{enumerate}
\def\labelenumi{\arabic{enumi})}
\setcounter{enumi}{6}
\tightlist
\item
  我们给单位圆 \(S_{1} \subset \mathbf{C}\) 赋以规范化哈尔测度。设 E 是复希尔伯特空间 \(\mathrm{L}^{2}(\mathbf{S}_{1})\) 。对 \(n \in \mathbf{Z}\) ，设 \(e_{n}\) 是 E 的元素 \(z \mapsto z^{n}\) 。设 F ⊂ E 是由诸 \(e_{n}\) （其中 \(n \geqslant 0\) ）生成的闭子空间，并设 \(p: E \to F\) 是正交投影。对 \(f \in L^{\infty}(\mathbf{S}_{1})\) ，我们在 E（相应地，F）上定义线性映射 \(M_{f}\)（相应地，\(T_{f}\)）如下：\(\mathrm{M}_{f}(g) = fg\) （\(g \in E\) ），且 \(T_{f} = p \circ M_{f}\) 。我们称 \(T_{f}\) 为以 f 为符号的特普利茨算子。(a) 线性映射 \(M_{f}\) 与 \(T_{f}\) 连续。\(M_{f}\) 的谱是 f 的傅里叶系数 \(c_{n} = \langle f | e_{n} \rangle\) 所成之集的闭包。自同态 \(M_{f}\) 是弗雷德霍姆自同态，当且仅当它是里斯自同态。
\end{enumerate}

\begin{enumerate}
\def\labelenumi{\alph{enumi})}
\setcounter{enumi}{1}
\item
  自同态 \(T_{e_{n}}\) （其中 \(n \in Z\) ）都是弗雷德霍姆自同态，且 \(\text{ind}(\mathrm{T}_{e_{n}}) = -n\) 对所有 \(n \in Z\) 成立。
\item
  \(\mathrm{T}_{\overline{f}} = \mathrm{T}_f^*\) 对每个函数 \(f\in \mathrm{L}^{\infty}(\mathbf{S}_1)\) 都成立。
\item
  若 \(T_{f}\) 是同构，则 \(M_{f}\) 是同构，且 f 在 \(\mathrm{L}^{\infty}(\mathbf{S}_{1})\) 中可逆。由此推出 \(M_{f}\) 的谱包含于 \(T_{f}\) 的谱，且 \(\|T_{f}\|=\|f\|\) 。
\item
  设 \(f \in \mathcal{C}(\mathbf{S}_1)\) 是连续函数，\(g \in \mathrm{L}^{\infty}(\mathbf{S}_1)\) 。自同态 \(\mathrm{T}_f\mathrm{T}_g - \mathrm{T}_{fg}\) 与 \(\mathrm{T}_g\mathrm{T}_f - \mathrm{T}_{gf}\) 是紧算子。（先考虑 \(f = e_n\) 的情形。）
\item
  设 \(\mathcal{C}alk(\mathrm{F})\) 是 F 的卡尔金代数，\(\pi\colon\mathcal{L}(\mathrm{F})\to\mathcal{C}alk(\mathrm{F})\) 是典范投影。证明 \(f\mapsto\pi(\mathrm{T}_{f})\) 是从 \(\mathcal{C}(\mathbf{S}_{1})\) 到 \(\mathcal{C}alk(\mathrm{F})\) 中的对合代数的单态射。（为证单性，可利用 I, p. 31 的命题 1。）
\end{enumerate}

\(g)\) 设 \(f \in \mathcal{C}(\mathbf{S}_1)\) 。则 \(\mathrm{T}_f\) 是弗雷德霍姆自同态，当且仅当 \(f\) 在 \(\mathbf{S}_1\) 上不取零值。（利用 I, p. 106 的命题 5。）

\(h)\) 设 \(f \in \mathcal{C}(\mathbf{S}_1)\) 且 \(f\) 在 \(\mathbf{S}_1\) 上不取零值。\(\mathrm{T}_f\) 的指标等于 \(n \in \mathbf{Z}\) ，当且仅当道路 \(t \mapsto f(e^{2i\pi t})\) （在 \(\mathbf{C} - \{0\}\) 中）同伦于道路 \(t \mapsto e_n(t)\) （参见 TA, III, p. 229, 定义 1）。

\begin{enumerate}
\def\labelenumi{\arabic{enumi})}
\setcounter{enumi}{7}
\tightlist
\item
  我们沿用上一习题的记号。
\end{enumerate}

\begin{enumerate}
\def\labelenumi{\alph{enumi})}
\item
  用 \(\mathrm{T}(\mathbf{S}_{1})\) 表示 \(\mathcal{L}(\mathrm{F})\) 中由特普利茨算子 \(T_{e_{1}}\) 与 \(T_{e_{-1}}\) 生成的闭子代数。证明 \(\mathrm{T}(\mathbf{S}_{1})\) 是 \(\mathcal{L}(\mathrm{F})\) 的星子代数。
\item
  设 \(f \in \mathcal{C}(\mathbf{S}_1)\) 。证明 \(T_f\) 属于 \(T(\mathbf{S}_1)\) 。
\item
  \(\mathrm{T}(\mathbf{S}_{1})\) 在 F 上的表示是不可约的；星代数 \(\mathrm{T}(\mathbf{S}_{1})\) 包含 \(\mathcal{L}^{\mathrm{c}}(\mathrm{F})\) 。
\end{enumerate}

\(d)\) \(u \in \mathrm{T}(\mathbf{S}_1)\) 当且仅当存在 \(f \in \mathcal{C}(\mathbf{S}_1)\) 与 \(h \in \mathcal{L}^c(\mathrm{F})\) 使得 \(u = \mathrm{T}_f + h\) 。（证明形如 \(\mathrm{T}_f + h\) 的自同态所成之集在 \(\mathcal{L}(\mathrm{F})\) 中构成一个闭星代数。）

\begin{enumerate}
\def\labelenumi{\alph{enumi})}
\setcounter{enumi}{4}
\tightlist
\item
  存在对合代数态射的一个正合列
\end{enumerate}

\[
0 \longrightarrow \mathcal {L} ^ {\mathrm{c}} (\mathrm{F}) \longrightarrow \mathrm{T} (\mathbf {S} _ {1}) \stackrel {\varrho} {\longrightarrow} \mathcal {C} (\mathbf {S} _ {1}) \longrightarrow 0
\]

使得 \(\varrho(\mathrm{T}_f) = f\) 对每个函数 \(f \in \mathcal{C}(\mathbf{S}_1)\) 成立。

\begin{enumerate}
\def\labelenumi{\arabic{enumi})}
\setcounter{enumi}{8}
\item
  设 B 是无穷维巴拿赫空间。用 E 表示 B 的强对偶，用 F 表示赋以拓扑 \(\sigma(\mathrm{B}^{\prime},\mathrm{B})\) 的 B 的对偶。证明：存在从 E 到 F 中的连续双射线性映射 u，以及从 E 到 F 中的紧线性映射 h，使得 \(u + h\) 有无穷维的余核。
\item
  设 E 是分离的局部凸空间，u 与 h 是 \(\mathcal{L}(\mathrm{E})\) 中交换的元素。假设存在整数 n 使得 \(h^{n}\) 是紧自同态。
\end{enumerate}

\begin{enumerate}
\def\labelenumi{\alph{enumi})}
\item
  若 u 是弗雷德霍姆自同态，则 u-h 是弗雷德霍姆自同态。
\item
  若 \(u\) 是里斯自同态，则 \(u - h\) 是里斯自同态。
\end{enumerate}

\begin{enumerate}
\def\labelenumi{\arabic{enumi})}
\setcounter{enumi}{10}
\tightlist
\item
  设 \(\mathrm{E} = \ell^2 (\mathbf{N})\) ，\((e_n)\) 是 E 的典范基。我们定义 \(u\in \mathcal{L}(\mathrm{E})\) 为满足 \(u(e_i) = e_{i + 1}\) （\(i\in \mathbf{N}\) ）的自同态。设 \(\mathrm{F} = \mathrm{E}\oplus \mathrm{E}\) ，\(v = u\oplus (u^{*} + 2\cdot 1_{\mathrm{E}})\in \mathcal{L}(\mathrm{F})\) 。设 \(\overline{v} = \pi (v)\) 是 \(v\) 在卡尔金代数 \(\mathcal{Calk}(\mathrm{F})\) 中的像。
\end{enumerate}

\begin{enumerate}
\def\labelenumi{\alph{enumi})}
\item
  自同态 \(v\) 是指标为 -1 的弗雷德霍姆自同态，而 \(v - 2 \cdot 1_{\mathrm{E}}\) 是指标为 1 的弗雷德霍姆自同态。
\item
  设 \(p = \mathrm{X}(\mathrm{X} - 2) \in \mathbf{C}[\mathrm{X}]\) ，S 是 \(\overline{v}\) 的指数谱（I, p. 173 的习题 5）。我们有 \(0 \in p(\mathrm{S})\) ，但 0 不属于 \(p(\overline{v})\) 的指数谱。（利用习题 3。）
\end{enumerate}

\exercisesection{}

\begin{enumerate}
\def\labelenumi{\arabic{enumi})}
\tightlist
\item
  设 A 是 \(\mathbf{C}\) 的紧子集，它等于其内部 V 的闭包。我们置 \(\mathcal{C}(\mathrm{A}) = \mathcal{C}(\mathrm{A};\mathbf{C})\) 。设 \(B(A)\) 是 \(\mathcal{C}(\mathrm{A})\) 中由在 V 上的限制为全纯的函数 \(f\in \mathcal{C}(\mathrm{A})\) 所组成的向量子空间；
\end{enumerate}

它是 \(\mathcal{C}(\mathrm{A})\) 的一个巴拿赫子代数。设 \(u \in \mathcal{L}(\mathrm{B}(\mathrm{A}))\) 是满足 \(u(f)(z) = zf(z)\) （对所有 \(z \in \mathrm{A}\) 及每个 \(f \in \mathrm{B}(\mathrm{A})\) ）的线性映射。

\begin{enumerate}
\def\labelenumi{\alph{enumi})}
\item
  \(u\) 的谱等于 A。
\item
  对每个 \(z \in V\) ，B(A) 的自同态 u - z 是指标为 1 的弗雷德霍姆自同态。
\item
  对每个 \(z \in A - V\) ，B(A) 的自同态 u - z 不是弗雷德霍姆自同态。
\item
  u 的稳定谱是 A 的边界，u 的敏感谱是空集，而 u 的本质谱是 A。
\item
  我们假设 A 在 \(\mathbf{C}\) 中的余集是连通的。对每个 \(h \in \mathcal{L}^{c}(\mathrm{B}(\mathrm{A}))\) ，\(u + h\) 的本质谱等于 A。
\end{enumerate}

\begin{enumerate}
\def\labelenumi{\arabic{enumi})}
\setcounter{enumi}{1}
\tightlist
\item
  我们保留上一习题的记号。假设 A 包含单位圆 \(U\) （\(\mathbf{C}\) 中的），且 \(0 \notin A\) 。用 \(\varphi_{A}\) 表示 A 的特征函数。
\end{enumerate}

\begin{enumerate}
\def\labelenumi{\alph{enumi})}
\tightlist
\item
  映射
\end{enumerate}

\[
\ell \colon f \mapsto \frac {1}{2 i \pi} \int_ {\mathbf {U}} f (z) d z
\]

是 B(A) 上的一个连续线性形式。

\begin{enumerate}
\def\labelenumi{\alph{enumi})}
\setcounter{enumi}{1}
\tightlist
\item
  设 \(h\) 是 B(A) 的有限秩自同态，满足 \(h(f) = \ell(f)\varphi_{\mathrm{A}}\) 。\(u - h\) 的本质谱异于 \(u\) 的本质谱。（证明它包含 0。）
\end{enumerate}

\begin{enumerate}
\def\labelenumi{\arabic{enumi})}
\setcounter{enumi}{2}
\tightlist
\item
  设 E 是复希尔伯特空间。设 U 是 \(\mathbf{C}\) 的连通开子集。设 f 是从 U 到空间 \(\mathcal{L}^{c}(\mathrm{E})\) 中的全纯函数。假设存在 \(z \in U\) 使得 1 属于 \(f(z)\) 的预解集。
\end{enumerate}

\begin{enumerate}
\def\labelenumi{\alph{enumi})}
\item
  由 \(z \in U\) （满足 \(1 \in \operatorname{Sp}(f(z))\) ）组成的集 D 是 U 的一个离散子集。
\item
  从 U-D 到 \(\mathcal{L}(\mathrm{E})\) 中的由 \(z\mapsto \mathrm{R}(f(z),1)\) 定义的映射是全纯的。
\end{enumerate}

我们固定一点 \(z_0 \in \mathrm{D}\) 。

\begin{enumerate}
\def\labelenumi{\alph{enumi})}
\setcounter{enumi}{2}
\item
  存在整数 \(k \geqslant 1\) ，使得从 U-D 到 \(\mathcal{L}(\mathrm{E})\) 中的由 \(z \mapsto (z - z_{0})^{k}\mathrm{R}(f(z), 1)\) 定义的映射可延拓为 \((\mathrm{U}-\mathrm{D}) \cup \{z_{0}\}\) 上的一个全纯映射。
\item
  存在 E 的自同态的唯一序列 \((u_{n})_{n\in\mathbf{N}}\) （当 n 充分大时 \(u_{n}=0\) ），以及从 \((\mathrm{U}-\mathrm{D})\cup\{z_{0}\}\) 到 \(\mathcal{L}(\mathrm{E})\) 中的唯一全纯函数 S，使得对所有 \(z\in U-D\) 有
\end{enumerate}

\[
\mathrm{R} (f (z), 1) = \sum_ {j = - k} ^ {- 1} (z - z _ {0}) ^ {j} u _ {j} + \mathrm{S} (z).
\]

\begin{enumerate}
\def\labelenumi{\alph{enumi})}
\setcounter{enumi}{4}
\tightlist
\item
  自同态 \(u_{-1}\) 是 \(f(z_0)\) 相对于 \(\{1\}\) （\(f(z_0)\) 的谱的开且闭子集）的谱投影。
\end{enumerate}

\begin{enumerate}
\def\labelenumi{\arabic{enumi})}
\setcounter{enumi}{3}
\tightlist
\item
  设 E 是可数型的巴拿赫空间。设 \(u \in \mathcal{L}(\mathrm{E})\) 。假设序列 \((\| u^k \|)_{k \in \mathbf{N}}\) 有界。
\end{enumerate}

对每个 \(\lambda\in\mathrm{Sp}(u)\cap\mathbf{U}\) ，设 \(x_{\lambda}\in\mathrm{E}\) 是 \(u\) 的属于 \(\lambda\) 的范数为 1 的特征向量 。

\begin{enumerate}
\def\labelenumi{\alph{enumi})}
\tightlist
\item
  对所有 \(\lambda_1\) 、\(\lambda_2\) （属于 \(\mathrm{Sp}(u) \cap \mathbf{U}\) ）以及每个整数 \(k \in \mathbf{N}\) ，我们有
\end{enumerate}

\[
\left| \lambda_ {1} ^ {k} - \lambda_ {2} ^ {k} \right| \leqslant \left(1 + \| u ^ {k} \|\right) \| x _ {\lambda_ {1}} - x _ {\lambda_ {2}} \|.
\]

\begin{enumerate}
\def\labelenumi{\alph{enumi})}
\setcounter{enumi}{1}
\item
  存在实数 c \textgreater{} 0 ，使得 \(\|x_{\lambda_{1}} - x_{\lambda_{2}}\| \geqslant c\) 对所有 \(\lambda_{1} \neq \lambda_{2}\) （\(\operatorname{Sp}(u) \cap \mathbf{U}\) 中的）成立。
\item
  集合 \(\operatorname{Sp}(u) \cap \mathbf{U}\) 是可数的（“贾米森定理”）。
\item
  若 E 不是可数型的，贾米森定理一般不成立。
\end{enumerate}

\begin{enumerate}
\def\labelenumi{\arabic{enumi})}
\setcounter{enumi}{4}
\tightlist
\item
  本习题的目标是在空间 E 为巴拿赫空间的情形给出洛莫诺索夫定理（III, p. 13 的推论 2）的另一个证明。
\end{enumerate}

设 E 是维数至少为 2 的复巴拿赫空间。设 h 是 E 的一个非零紧自同态，A 是 h 在 \(\mathcal{L}(\mathrm{E})\) 中的交换子。a) 下列条件之一成立：

\begin{enumerate}
\def\labelenumi{(\roman{enumi})}
\item
  存在 \(y \in E\) ，使得 Ay 是 E 中的一个非零且不稠密的子空间；
\item
  对 A 的每个内部非空的闭有界子集 C，存在 A 的有限子集 I，使得 \(h(\mathrm{C})\) 包含于诸集合 \(u^{-1}(\mathrm{C})\) （\(u \in I\) ）的并之中。
\end{enumerate}

\begin{enumerate}
\def\labelenumi{\alph{enumi})}
\setcounter{enumi}{1}
\tightlist
\item
  假设上一问题的条件 (ii) 成立。设 C 是 E 的一个内部非空的闭有界子集。对每个整数 \(m \geqslant 1\) ，存在 A 中的有限序列 \((u_{1}, \ldots, u_{m})\) 使得
\end{enumerate}

\[
u _ {1} \circ \dots \circ u _ {m} \circ h ^ {m} (x _ {0}) \in \mathrm{C}.
\]

（打乒乓球。）

\begin{enumerate}
\def\labelenumi{\alph{enumi})}
\setcounter{enumi}{2}
\item
  在同一假设下，证明 h 的谱半径非零（在相反的情形，考虑一个元素 \(x_{0} \in E\) ，满足 \(\|h(x_{0})\| > \|h\|\) ，以及以 \(x_{0}\) 为心、\(\|h\|\) 为半径的闭球 C；注意到 \(\|( \lambda h)^{n}\|\) 对每个 \(\lambda \in R\) 趋于 0，证明上一问题的结论将蕴含 0 属于 C 的闭包）。
\item
  对 E 的每个与 E 的某个非零紧自同态交换的自同态 \(u\) ，存在异于 E 的非零闭子空间 \(F \subset E\) 使得 \(u(F) \subset F\) 。
\end{enumerate}

\begin{enumerate}
\def\labelenumi{\arabic{enumi})}
\setcounter{enumi}{5}
\tightlist
\item
  设 A 是一个幺 Z-代数。假设 A 是自由 Z-模。设 B 是 A 的一个子集，它作为 Z-模是 A 的基。我们称对子 \((A, B)\) 是一个自然代数，如果 A 关于 B 的结构常数（A, III, p. 10）属于 N，且 \(1_{A}\) 在基 B 中的系数属于 N。此外，我们称 (A, B) 是幺的，如果 \(1_{A}\) 是 B 的元素。
\end{enumerate}

设 (A, B) 是一个自然代数。设 \(B_{0}\) 是由 \(b \in B\) （使 b 在 \(1_{A}\) 中的系数非零）组成的集。用 \(\tau_{A}\) 表示从 A 到 Z 的唯一的 Z-模同态，使得 \(\tau_{\mathrm{A}}(b) = 1\) （若 \(b \in B_{0}\) ），否则 \(\tau_{\mathrm{A}}(b) = 0\) 。

\begin{enumerate}
\def\labelenumi{\alph{enumi})}
\item
  \(\tau_{\mathrm{A}}(xy) = \tau_{\mathrm{A}}(yx)\) 对所有 \((x,y)\in \mathrm{A}^2\) 成立。
\item
  对所有 \(b \neq b'\) （在 \(B_{0}\) 中），有 \(b^2 = b\) 且 \(bb' = 0\) 。
\end{enumerate}

\begin{enumerate}
\def\labelenumi{\arabic{enumi})}
\setcounter{enumi}{6}
\tightlist
\item
  设 (A, B) 是一个自然代数。
\end{enumerate}

\begin{enumerate}
\def\labelenumi{\alph{enumi})}
\tightlist
\item
  至多存在一个对合 \(\iota\) （\(\mathcal{B}\) 的），使得从 A 到 A 的唯一 \(\mathbf{Z}\)-模同态（把 \(b\) 映为 \(\iota(b)\) ）是 A 的一个反自同构，且满足条件 \(\tau_{\mathrm{A}}(bb') = 1\) （若 \(b' = \iota(b)\) ），否则 \(\tau_{\mathrm{A}}(bb') = 0\) 。
\end{enumerate}

当这样一个对合 \(\iota\) 存在时，我们称 (A, \(\mathcal{B}\) ) 是一个自然基础代数。

我们现在假设 (A, B) 是一个自然基础代数。

\begin{enumerate}
\def\labelenumi{\alph{enumi})}
\setcounter{enumi}{1}
\tightlist
\item
  于是 \(\iota(b) = b\) 对所有 \(b \in B_{0}\) 成立，并有公式
\end{enumerate}

\[
1 _ {\mathrm{A}} = \sum_ {b \in \mathcal {B} _ {0}} b.
\]

\begin{enumerate}
\def\labelenumi{\alph{enumi})}
\setcounter{enumi}{2}
\tightlist
\item
  A 的结构常数 \(\gamma_{bb'}^{b''}\) 满足 \(\gamma_{bb'}^{\iota(b'')} = \tau_{\mathrm{A}}(bb'b'')\) （对 \((b, b', b'')\) ，\(\mathcal{B}^{3}\) 中的）。
\end{enumerate}

\begin{enumerate}
\def\labelenumi{\arabic{enumi})}
\setcounter{enumi}{7}
\tightlist
\item
  设 (A, B) 是一个自然基础代数，且 A 的秩有限。
\end{enumerate}

\begin{enumerate}
\def\labelenumi{\alph{enumi})}
\tightlist
\item
  对每个 \(x \in A\) ，元素
\end{enumerate}

\[
\sum_ {b \in \mathcal {B}} b x \iota (b)
\]

是中心元。

在本习题的其余部分，我们进一步假设 (A, \(\mathcal{B}\) ) 是幺的。

\begin{enumerate}
\def\labelenumi{\alph{enumi})}
\setcounter{enumi}{1}
\item
  对所有 \(b\) 与 \(b'\) （\(\mathcal{B}\) 中的），存在 \(b_1\) 与 \(b_2\) （\(\mathcal{B}\) 中的），使得 \(b'\) 在 \(bb_1\) 中以及在 \(b_2b\) 中的系数非零。
\item
  对每个 \(x \in B\) ，在基 B 中表示从 A 到 A 的映射 \(y \mapsto xy\) 的矩阵是一个实矩阵，其所有系数都是正的。用 pf 表示从 A 到 R 的唯一同态，使得 \(\operatorname{pf}(b)\) （当 \(b \in B\) 时）是这个矩阵的谱半径。
\end{enumerate}

\(d)\) 对每个 \(b \in \mathcal{B}\) ，实数 \(\alpha = \mathrm{pf}(b)\) 是一个代数整数，满足 \(\alpha \geqslant 1\) ；对每个共轭数 \(\alpha' \in \mathbf{C}\) （\(\alpha\) 的），有 \(|\alpha'| \leqslant \alpha\) 。

\begin{enumerate}
\def\labelenumi{\alph{enumi})}
\setcounter{enumi}{4}
\tightlist
\item
  设 \(a_0 \in \mathrm{A}\) 是 \(\mathcal{B}\) 的诸元素之和。表示映射 \(f: y \mapsto ya_0\) 的矩阵在基 \(\mathcal{B}\) 中的所有系数都严格为正。存在一个元素 \(r \in \mathbf{C} \otimes \mathrm{A}\) ，它在基 \((1 \otimes b)_{b \in \mathcal{B}}\) 中的系数都严格为正，使得 \(f_{\mathbf{C}}(r) = \varrho(f_{\mathbf{C}})r\) 。
\end{enumerate}

\(f)\) 对每个 \(x \in \mathrm{A}\) ，存在唯一的复数 \(\delta(x)\) ，使得 \(xr = \delta(x)r\) 在 \(\mathbf{C} \otimes \mathrm{A}\) 中成立。映射 \(\delta: \mathrm{A} \to \mathbf{C}\) 是 \(\mathrm{A}\) 的一个特征标。

\begin{enumerate}
\def\labelenumi{\alph{enumi})}
\setcounter{enumi}{6}
\tightlist
\item
  \(\delta(x) = \mathrm{pf}(x)\) 对所有 \(x \in \mathrm{A}\) 成立。特别地，映射 \(x \mapsto \mathrm{pf}(x)\) 是 A 的一个特征标。
\end{enumerate}

\(h)\) 映射 pf 是 A 的唯一特征标，使得 \(\mathrm{pf}(b) \geqslant 0\) 对所有 \(b \in \mathcal{B}\) 成立；此外，\(\mathrm{pf}(b) > 0\) 对所有 \(b \in \mathcal{B}\) 成立。

\begin{enumerate}
\def\labelenumi{\roman{enumi})}
\tightlist
\item
  对每个 \(y \in \mathrm{A}\) ，有 \(ry = \delta(y)r\) （在 \(\mathbf{C} \otimes \mathbf{A}\) 中）。
\end{enumerate}

\begin{enumerate}
\def\labelenumi{\alph{enumi})}
\setcounter{enumi}{9}
\tightlist
\item
  元素
\end{enumerate}

\[
r _ {0} = \sum_ {b \in \mathcal {B}} \mathrm{pf} (b) b
\]

具有问题 i) 所要求的性质。

\begin{enumerate}
\def\labelenumi{\arabic{enumi})}
\setcounter{enumi}{8}
\tightlist
\item
  \begin{enumerate}
  \def\labelenumii{\alph{enumii})}
  \tightlist
  \item
    设 \(n\) 是整数，\(M\) 是系数属于 \(\mathbf{N}\) 的一个方阵。假设 \(\varrho(M^t M) = \varrho(M)^2\) 且 \(|\varrho(M)| < 2\) 。则存在整数 \(m \geqslant 2\) 使得 \(\varrho(M) = 2\cos(\pi/m)\) 。（利用克罗内克定理，参见 A, V, p. 155, 习题 17。）
  \end{enumerate}
\end{enumerate}

\begin{enumerate}
\def\labelenumi{\alph{enumi})}
\setcounter{enumi}{1}
\item
  对每个整数 \(m \geqslant 2\) ，存在一个系数属于 N 的方阵 M，使得 \(\varrho(M) = 2 \cos(\pi/m)\) 。
\item
  设 (A, B) 是一个秩有限的幺自然基础代数。设 \(b \in B\) 。若 \(\operatorname{pf}(b) < 2\) ，则存在整数 \(m \geqslant 3\) 使得 \(\operatorname{pf}(b) = 2 \cos(\pi/m)\) 。
\end{enumerate}

习题 10 至 24 给出 III, p. 100 的推论 5 的另一个证明以及若干强化。其中将用到下列概念。

设 I 是有限集，n 是它的基数。

我们用 \(1_{\mathrm{I}}\) 表示 (I, I) 型的实单位矩阵。

- 一个 (I, I) 型的实系数方阵称为正的，如果它的所有系数都 \(\geqslant 0\) ；称为严格正的，如果它的所有系数都 \(>0\) 。

- (I,I) 型的实系数方阵 \((a_{ij})_{i,j\in \mathrm{I}}\) 称为可约的，如果存在把 I 分成两个非空子集 I' 与 I'\,' 的划分，使得 \(a_{ij} = 0\) 对所有 \((i,j)\in \mathrm{I}'\times \mathrm{I}''\) 成立；若它不是可约的，则称为不可约的。

- 我们把 \(\mathbf{R}^{\mathrm{I}}\) 的由典范基向量的某个子集生成的任一向量子空间称为 \(\mathbf{R}^{\mathrm{I}}\) 的坐标子空间。坐标子空间共有 \(2^{n}\) 个；对每个整数 \(k\in\{0,\ldots,n\}\) ，其中维数为 \(k\) 的坐标子空间有 \(\binom{n}{k}\) 个。

因此，(I, I) 型的实矩阵是不可约的，当且仅当它所稳定的坐标子空间只有 \(\{0\}\) 和 \(R^{I}\) 。

对所有 \(x\) 与 \(y\) （\(\mathbf{R}^{\mathrm{I}}\) 中的），我们记 \(x \preccurlyeq y\) （及 \(y \succcurlyeq x\) ），如果 \(y - x\) 属于 \(\mathbf{R}^{\mathrm{I}}\) 中坐标为正的元素所成的锥，亦即 \(x_i \leqslant y_i\) 对所有 \(i \in \mathrm{I}\) 成立。我们记 \(x \prec y\) （及 \(y \succ x\) ），如果 \(y - x\) 的所有坐标都严格为正。

\begin{enumerate}
\def\labelenumi{\arabic{enumi})}
\setcounter{enumi}{9}
\item
  设 \(A\) 是 (I, I) 型的正且不可约的实矩阵。对每个整数 \(n\) （满足 \(n \geqslant \operatorname{Card}(\mathrm{I})\) ），矩阵 \((1_{\mathrm{I}} + A)^{n-1}\) 是严格正的。（设 \(y \in \mathbf{R}_{+}^{n}\) ，\(z = (1_{\mathrm{I}} + A)y\) ；若 \(y \neq 0\)，证明由 \(i \in \mathrm{I}\) （使 \(z_{i} = 0\) ）所成的集严格包含于 \(\{i \in \mathrm{I} \mid y_{i} = 0\}\) 。）
\item
  设 \(A\) 是 (I, I) 型的正且不可约的实矩阵。用 \(\psi \in \mathbf{R}[\mathrm{X}]\) 表示 \(A\) 的极小多项式，并置 \(m = \deg(\psi)\) 。
\end{enumerate}

对每个整数 \(q\)，置 \(A^q = (a_{ij}^{(q)})_{i,j \in \mathrm{I}}\) 。

对每个 \((i,j)\in\mathrm{I}^{2}\) ，存在整数 q 使得 \(1\leqslant q\leqslant m\) 且 \(a_{ij}^{(q)}>0\) 。（设 r 是 \((1+\mathrm{X})^{n-1}\) 除以 \(\psi\) 的欧几里得除法的余式；考察矩阵 \(r(A)\) 并注意它是严格正的。）

\begin{enumerate}
\def\labelenumi{\arabic{enumi})}
\setcounter{enumi}{11}
\tightlist
\item
  设 A 是 (I, I) 型的不可约且正的实矩阵。对每个非零的 \(x \in R^{I}\) ，我们置
\end{enumerate}

\[
r_{x} = \inf_{\substack{i\in \mathrm{I}\\ x_{i}\neq 0}}\frac{(Ax)_{i}}{x_{i}},
\]

其中 \((Ax)_i = \sum_{k\in I}a_{ik}x_k\) 是 \(Ax\) 的以 \(i\) 为指标的坐标。

\begin{enumerate}
\def\labelenumi{\alph{enumi})}
\item
  若 \(x \succcurlyeq 0\) 且 \(x \neq 0\) ，则 \(r_x\) 是最大的实数 \(\varrho\) ，使得 \(\varrho x \preccurlyeq Ax\) 。
\item
  假设 \(x \succcurlyeq 0\) ，并置 \(y = (1_{\mathrm{I}} + A)x\) 。则有 \(r_x \leqslant r_y\) 。
\item
  限制在 \((\mathbf{R}_{+}^{*})^{\mathrm{I}}\) 上时，函数 \(x \mapsto r_{x}\) 是连续的。
\end{enumerate}

\(d)\) 由此推出存在 \(z \in \mathbf{R}^{\mathrm{I}}\) 满足 \(z \succcurlyeq 0\) 且 \(r_z = \sup_{x \succcurlyeq 0} r_x\) （利用紧集 \(\mathrm{M} = \{x \in \mathbf{R}_{+}^{\mathrm{I}} \mid \sum_{i \in \mathrm{I}} x_i = 1\}\) 以及 \((1_{\mathrm{I}} + A)^{n-1}(\mathrm{M})\) ，再用习题 10。）

我们置 \(r = r_{z}\) 。\(R^{I}\) 中满足 \(z \succcurlyeq 0\) 且 \(r_{z} = r\) 的元素 z 称为极端元素。

\begin{enumerate}
\def\labelenumi{\alph{enumi})}
\setcounter{enumi}{4}
\item
  我们有 r \textgreater{} 0 。
\item
  若 z 是极端向量，则 Az = rz。（若 \(u = Az - rz \succcurlyeq 0\) 非零，则由习题 10，\((1_{\mathrm{I}} + A)^{n-1}u\) 将为 \(\succ 0\) ，且对 \(x = (1_{\mathrm{I}} + A)^{n-1}z\) 将有 \(Ax \succ rx\) 。）
\item
  每个极端元素都 \(\succ 0\) 。
\end{enumerate}

对每个向量 \(y\) （在 \(\mathbf{C}^{\mathrm{I}}\) 中），用 \(y^{+}\) 表示向量 \((|y_{i}|)_{i\in \mathrm{I}}\) ；于是有 \(y^{+}\in \mathbf{R}^{\mathrm{I}}\) 且 \(y^{+}\succcurlyeq 0\) 。

\begin{enumerate}
\def\labelenumi{\alph{enumi})}
\setcounter{enumi}{7}
\item
  设 \(\alpha \in \mathbf{C}\) 是把 \(A\) 视为复矩阵时的一个特征值，\(y\) 是 \(\mathbf{C}^{\mathrm{I}}\) 中相应的一个特征向量。我们有 \(|\alpha|y^{+} \preccurlyeq Ay^{+}\) 且 \(|\alpha| \leqslant r\) 。
\item
  推出：若 \(\alpha = r\) ，则 y 的所有坐标都非零，进而 A 的属于特征值 r 的特征空间（在 \(C^{I}\) 中）的维数为 1。
\end{enumerate}

用 \(\Delta(\mathrm{X})\) 表示 \(\mathrm{X} \cdot 1_{\mathrm{I}} - A\) 的行列式，用 \(B(\mathrm{X})\) 表示 \(\mathrm{X} \cdot 1_{\mathrm{I}} - A\) 的余子式矩阵的转置。我们记 \(B(\mathrm{X}) = (B_{ij}(\mathrm{X}))_{i,j \in \mathrm{I}}\) 。于是有 \(B(\mathrm{X})(\mathrm{X} \cdot 1_{\mathrm{I}} - A) = (\mathrm{X} \cdot 1_{\mathrm{I}} - A)B(\mathrm{X}) = \Delta(\mathrm{X}) \cdot 1_{\mathrm{E}}\) （A, III, p. 99, 命题 12）。

\begin{enumerate}
\def\labelenumi{\alph{enumi})}
\setcounter{enumi}{9}
\tightlist
\item
  对每个 \(i \in I\) ，\(B(r)\) 的以 i 为指标的列非零。（若不然，构造 A 的属于特征值 r 的一个特征向量，其第 i 个系数为零，这与 f) 矛盾。）
\end{enumerate}

\(k\) ) \(B(r)\) 的所有系数都有相同的符号。（首先利用关系式 \((r1_{\mathrm{I}} - A)B(r) = 0\) 证明：对所有 \(i\in \mathrm{I}\) ，以 \(i\) 为指标的 \(B(r)\) 的列的所有系数非零且同号，然后再对转置进行论证。）

\(l)\) \(\Delta'(r) > 0\) 且 \(B_{ij}(r) > 0\) 对所有 \((i,j) \in \mathrm{I}^2\) 成立（利用公式 \(\Delta'(X) = \sum_i B_{ii}(X)\) ，并注意 \(\Delta'\) 在区间 \([r, +\infty]\) 上不变号。）

\begin{enumerate}
\def\labelenumi{\alph{enumi})}
\setcounter{enumi}{12}
\tightlist
\item
  由此得出 r 是 A 的特征多项式的单根。
\end{enumerate}

\begin{enumerate}
\def\labelenumi{\arabic{enumi})}
\setcounter{enumi}{12}
\tightlist
\item
  设 A（相应地，C）是 (I, I) 型的实（相应地，复）方阵。假设 A 是正的且不可约，并且 \(C^{+} \preccurlyeq A\) ，这里与习题 12 一样，用 \(C^{+}\) 表示其系数为 C 的系数的绝对值的矩阵，而记号 \(C^{+} \preccurlyeq A\) 表示 \(A - C^{+}\) 的系数都是正的。
\end{enumerate}

我们沿用上一习题的记号。

\begin{enumerate}
\def\labelenumi{\alph{enumi})}
\item
  设 \(\gamma \in \mathbf{C}\) 是 \(C\) 的一个特征值，\(y\) 是 \(\mathbf{C}^{\mathrm{I}}\) 中相应的一个特征向量。我们有 \(|\gamma|y^{+} \preccurlyeq C^{+}y^{+} \preccurlyeq Ay^{+}\) 且 \(|\gamma| \leqslant r_{y^{+}} \leqslant r\) 。
\item
  我们进一步假设 \(|\gamma| = r\) 。则 \(y^{+}\) 是 A 的一个极端向量（参见习题 12 的问题 f)），且有关系式 \(Ay^{+} = C^{+}y^{+} = ry^{+}\) 与 \(A = C^{+}\) 。
\end{enumerate}

设 \(u\)（相应地，\(u_j\) ，\(j \in I\) ）是使 \(\gamma = |\gamma|u\) （相应地，\(y_j = |y_j|u_j\) ，\(j \in I\) ）成立的唯一复数。用 \(D\) 表示 (I, I) 型的复方阵，其对角系数为 \((u_j)_{j \in I}\) ，其余系数为零，于是 \(y = Dy^+\)。设 \(F = u^{-1}D^{-1}CD\) 。

\begin{enumerate}
\def\labelenumi{\alph{enumi})}
\setcounter{enumi}{2}
\item
  有 \(Fy^{+} = ry^{+}\) 且 \(F^{+}y^{+} = Fy^{+}\) （注意 \(F^{+} = C^{+}\) 。）
\item
  有 \(F = F^{+}\) 且 \(C = uDAD^{-1}\) 。
\end{enumerate}

\begin{enumerate}
\def\labelenumi{\arabic{enumi})}
\setcounter{enumi}{13}
\tightlist
\item
  设 A 是 (I, I) 型的正且不可约的实方阵。我们沿用习题 12 的记号。
\end{enumerate}

设 \((\lambda_{h})_{h\in\mathrm{H}}\) 是 A 的模为 r 的复特征值按重数重复而得的族；置 \(v_{h}=\lambda_{h}/r\) （对所有 \(h\in H\) ）。

设 \(h \in \mathrm{H}\) ；把习题 13 应用于 \(C = A\) 与 \(\gamma = \lambda_h\) ，存在矩阵 \(D_{(h)}\) 使得 \(A = v_h D_{(h)} AD_{(h)}^{-1}\) 且 \(D_{(h)}^+ = 1_{\mathrm{I}}\) 。

\begin{enumerate}
\def\labelenumi{\alph{enumi})}
\item
  设 \(z \succ 0\) 是 A 的一个极端向量。对 \(h \in H\) ，向量 \(y_{(h)} = D_{(h)}z\) 是 A 的属于特征值 \(\lambda_{h}\) 的特征向量。
\item
  证明：对 \(h \in H\) ，A 的特征值 \(\lambda_{h}\) 是单重的。由此推出矩阵 \(D_{(h)}\) 与向量 \(y_{(h)}\) 在相差一个纯量因子的意义下唯一确定。
\item
  证明：对 \(H^{2}\) 中所有 h 与 k，有 \(A = u_{h}u_{k}^{-1}D_{(h)}D_{(k)}^{-1}AD_{(k)}D_{(h)}^{-1}\) ，从而向量 \(D_{(h)}D_{(k)}^{-1}z\) 是 A 的属于特征值 \(u_{h}u_{k}^{-1}r\) 的特征向量。
\item
  由此推出 \(\{u_h\}_{h\in \mathrm{H}}\) 是 \(\mathbf{C}^*\) 的一个子群，且存在非零复数的一个族 \((\mu_h)_{h\in \mathrm{H}}\) ，使得 \(\{\mu_h^{-1}D_{(h)}\}_{h\in \mathrm{H}}\) 是 (I,I) 型复系数可逆矩阵所成的群的一个子群，并且这两个子群都是阶为 Card(H) 的循环群。（选取一个指标 \(i\in \mathrm{I}\) ，并把 \(\mathrm{D}_{(h)}\) 规范化，使其行指标与列指标均为 \(i\) 的系数等于 1。）
\end{enumerate}

在本习题的其余部分，我们假设 H = Z/dZ，且存在 \(\omega \in C^{*}\) 使得 \(\omega^{d} = 1\) 且 \(u_{h} = \omega^{h}\) （\(h \in H\) ）。置 \(D = D_{(1)}\) ，则 \(D^{d} = 1_{I}\) 且 \(D_{(h)} = D^{h}\) 对所有 \(h \in H\) 成立。

\begin{enumerate}
\def\labelenumi{\alph{enumi})}
\setcounter{enumi}{4}
\tightlist
\item
  由此推出 \(A = \omega DAD^{-1}\) 。
\end{enumerate}

我们考虑 I 的划分 \((\mathrm{I}_{h})_{h\in\mathrm{H}}\) ，使得对 \(i\in I\) ，D 的以 \((i,i)\) 为指标的系数等于 \(\omega^{h}\) 当且仅当 \(i\in I_{h}\) 。对每个 h，置 \(n_{h}=\mathrm{Card}(I_{h})\) 。矩阵 D 写成“分块”形式

\[
D = \left( \begin{array}{c c c c} 1 _ {n _ {0}} & & & \\ & \omega 1 _ {n _ {1}} & & \\ & & \ddots & \\ & & & \omega^ {d - 1} 1 _ {n _ {d - 1}} \end{array} \right)
\]

我们考虑 \(A\) 的相应的分解：

\[
A = \left( \begin{array}{c c c c} A _ {0, 0} & A _ {0, 1} & \dots & A _ {0, d - 1} \\ A _ {1, 0} & A _ {1, 1} & \dots & A _ {1, d - 1} \\ \vdots & \vdots & \ddots & \vdots \\ A _ {d - 1, 0} & A _ {d - 1, 1} & \dots & A _ {d - 1, d - 1} \end{array} \right),
\]

使得对 \(h, k \in \mathrm{H}\) ，\(A_{h,k}\) 是 \((\mathrm{I}_h, \mathrm{I}_k)\) 型的矩阵。

\(f)\) 对所有 \((h,k)\in\mathrm{H}\) ，有 \(\omega A_{h,k}=\omega^{k-h}A_{h,k}\) 。由此推出 \(\mathrm{A}_{h,k}=0\) （若 \(k-h\neq1\) ）。

\begin{enumerate}
\def\labelenumi{\arabic{enumi})}
\setcounter{enumi}{14}
\tightlist
\item
  我们沿用习题 12 的假设与记号。设 \(\Delta_1(\mathrm{X})\) 是 \(B(\mathrm{X})\) 的诸系数的最大公因子；假设多项式 \(\Delta_1(\mathrm{X})\) 是首一的。
\end{enumerate}

\begin{enumerate}
\def\labelenumi{\alph{enumi})}
\item
  多项式 \(\Delta_1(X)\) 整除 \(\Delta(X)\) ，且 \(\Delta_1(r) \neq 0\) 。
\item
  由此推出 \(\Delta_{1}(X)\) 的实根都严格小于 r。从而得出 \(\Delta_{1}(r)>0\) 。
\item
  置 \(C(\mathrm{X}) = \frac{1}{\Delta_1(\mathrm{X})} B(\mathrm{X})\) 。这是一个系数属于 \(\mathbf{R}[\mathrm{X}]\) 的矩阵。证明 \(C(r) \succ 0\) 。
\end{enumerate}

\begin{enumerate}
\def\labelenumi{\arabic{enumi})}
\setcounter{enumi}{15}
\tightlist
\item
  设 A 是 (I, I) 型的正且不可约的实方阵。对所有 \(j \in I\) 置 \(s_{j} = \sum_{k \in I} a_{jk}\) ，
\end{enumerate}

\[
s = \inf _ {j \in \mathrm{I}} s _ {j}, \quad \text {and} \quad \mathrm{S} = \sup _ {j \in \mathrm{I}} s _ {j}.
\]

我们也采用习题 12 的记号 \(r\) 与 \(B(\mathbf{X})\) 。

\begin{enumerate}
\def\labelenumi{\alph{enumi})}
\item
  有 \(\sum_{j\in \mathrm{I}}(r - s_j)B_{kj}(r) = 0\)
\item
  由此推出 \(s \leqslant r \leqslant S\) ，且这两个不等式中任何一个取等号都蕴含 s = S。
\end{enumerate}

\begin{enumerate}
\def\labelenumi{\arabic{enumi})}
\setcounter{enumi}{16}
\item
  设 \(A\) 是 (I, I) 型的正且不可约的实方阵。设 \(y \in \mathbf{R}^{\mathrm{I}}\) 是 \(A\) 的属于特征值 \(\alpha\) 的特征向量，且满足 \(y \succcurlyeq 0\) 。则 \(\alpha = r\) （沿用习题 12 的记号）且 \(y \succ 0\) 。（考虑 \(u\) （\(^{t}A\) 的一个极端向量），并计算 \(\langle y|^{t}Au\rangle = \langle Ay|u\rangle\) 。）
\item
  设 A 是 (I, I) 型的正且不可约的实方阵。设 \(\varrho\) 是 A 的谱半径。对每个 \(x \in R^{I}\) （满足 \(y \succcurlyeq 0\) 且 \(x \neq 0\) ），我们置
\end{enumerate}

\[
\varrho^ {x} = \sup _ {i \in \mathrm{I}} \frac {(A x) _ {i}}{x _ {i}},
\]

（约定 \(\varrho^x = +\infty\) ，如果存在指标 \(i\) 使 \(x_i = 0\) 且 \((Ax)_i \neq 0\) ）。设

\[
\varrho^{\prime} = \inf_{\substack{x\succ 0\\ x\neq 0}}\varrho^{x}.
\]

沿用与习题 12 类似的途径，证明 \(\varrho'\) 是 \(A\) 的一个特征值且 \(\varrho' = \varrho\) 。

\begin{enumerate}
\def\labelenumi{\arabic{enumi})}
\setcounter{enumi}{18}
\tightlist
\item
  设 A 是 (I, I) 型的正且不可约的实方阵；\(\varrho\) 是它的谱半径。设 \(z \in R^{I}\) 满足 \(z \succcurlyeq 0\) 且 \(z \neq 0\) 。
\end{enumerate}

若 \(\varrho z \preccurlyeq Az\) （相应地，\(\varrho z \succcurlyeq Az\) ），则 \(\varrho z = Az\) 且 \(z \succ 0\) 。

\begin{enumerate}
\def\labelenumi{\arabic{enumi})}
\setcounter{enumi}{19}
\tightlist
\item
  设 A 是 (I, I) 型的正但未必不可约的实方阵；\(\varrho\) 是 A 的谱半径。我们保留前面诸习题的记号 \(B(\mathrm{X})\) 与 \(C(\mathrm{X})\) 。
\end{enumerate}

\begin{enumerate}
\def\labelenumi{\alph{enumi})}
\item
  实数 \(\varrho\) 是 \(A\) 的特征值，且存在至少一个特征向量 \(\succcurlyeq 0\) 相伴于 \(\varrho\) 。
\item
  对每个实数 \(\lambda \geqslant r\) ，矩阵 \(B(\lambda)\) 与 \(B'(\lambda)\) 都是正的。（对第二个矩阵，可对 I 的基数用归纳法。）
\item
  矩阵 \(C(\lambda)\) 对所有 \(\lambda \geqslant r\) 都是正的。
\item
  若 A 是不可约的，则对每个实数 \(\lambda > r\) ，矩阵 \(B(\lambda)\) 、\(C(\lambda)\) 与 \((\lambda1_{\mathrm{I}} - A)^{-1}\) 都是严格正的。
\item
  设 J 是 I 的异于 I 的子集。设 \(\varrho_{\mathrm{J}}\) 是子矩阵 \((a_{i,j})_{i,j\in \mathrm{J}}\) 的谱半径。我们有 \(\varrho_{\mathrm{J}}\leqslant \varrho\) ，且当 \(A\) 不可约时该不等式是严格的。反之，若 \(A\) 是可约的，则存在异于 I 的子集 \(\mathrm{J}\subset \mathrm{I}\) 使得 \(\varrho_{\mathrm{J}} = \varrho\) 。
\item
  设 J 是 I 的异于 I 的子集，且子矩阵 \((a_{i,j})_{i,j\in\mathrm{J}}\) 的谱半径等于 \(\varrho\) 。对每个满足 \(J\subset K\subset I\) 的子集 K，矩阵 \((a_{i,j})_{i,j\in\mathrm{K}}\) 的谱半径等于 \(\varrho\) 。由此推出 \(A\succ0\) 当且仅当 \(B(r)\) 的所有对角系数都严格为正。
\end{enumerate}

\(g)\) 对每个实数 \(\lambda >\varrho\) 以及 I 的每个满足 \(\mathrm{J}\neq \mathrm{I}\) 的子集 J，矩阵 \((\lambda \delta_{i,j} - a_{i,j})_{i,j\in \mathrm{J}}\) 的行列式严格为正。

\begin{enumerate}
\def\labelenumi{\alph{enumi})}
\setcounter{enumi}{7}
\tightlist
\item
  设 \(\lambda\) 是实数；假设 \((\lambda \delta_{i,j} - a_{i,j})_{i,j \in \mathbf{J}}\) 的行列式对每个子集 \(\mathbf{J} \subset \mathbf{I}\) （异于 \(\mathbf{I}\) ）都严格为正。则有 \(\lambda > \varrho\) 。
\end{enumerate}

在下面的问题中，我们假设 \(\mathrm{I}\) 是区间 \([1,n]\) （在 \(\mathbf{N}\) 中）。i) 设 \((g_{i,j})_{1\leqslant i,j\leqslant n}\) 是一个实矩阵，使得 \(g_{i,j}\leqslant0\) 对所有 \(i\neq j\) 成立，且矩阵 \((g_{i,j})_{1\leqslant i,j\leqslant m}\) 的行列式对每个整数 \(m\) （满足 \(1\leqslant m\leqslant n\) ）都严格为正。则对 I 的每个子集 J 有 \(\det(g_{i,j})_{i,j\in\mathrm{J}}>0\) 。

\begin{enumerate}
\def\labelenumi{\alph{enumi})}
\setcounter{enumi}{9}
\item
  设 \(\lambda\) 是实数；假设对每个整数 \(m\) （满足 \(1 \leqslant m \leqslant n\) ），行列式 \(\det(\lambda \delta_{i,j} - a_{i,j})_{1 \leqslant i,j \leqslant m}\) 严格为正。则 \(\lambda > \varrho\) 。
\item
  设 \(C = (c_{i,j})_{1 \leqslant i,j \leqslant n}\) 是一个实矩阵，使得 \(c_{i,j} \geqslant 0\) 对所有 \(i \neq j\) 成立。若 \((-1)^{k} \det(c_{i,j})_{1 \leqslant i,j \leqslant m}\) 对每个满足 \(1 \leqslant m \leqslant n\) 的整数 m 都严格为正，则 C（视为复矩阵）的特征值的实部都是负的。
\end{enumerate}

\begin{enumerate}
\def\labelenumi{\arabic{enumi})}
\setcounter{enumi}{20}
\tightlist
\item
  设 A 是 (I, I) 型的正的实方阵；\(\varrho\) 是它的谱半径。
\end{enumerate}

\begin{enumerate}
\def\labelenumi{\alph{enumi})}
\tightlist
\item
  证明：存在整数 p 与 q 满足 \(0 \leqslant q \leqslant p\) ，以及 I 的一个分成非空部分的划分 \((\mathrm{I}_{s})_{1 \leqslant s \leqslant p}\) ，使得相应的分块分解 \(A = (A_{s,t})_{1 \leqslant s, t \leqslant p}\) 具有下列性质：
\end{enumerate}

- 矩阵 \(A_{s,s}\) 都是不可约的；

- 我们有 \(A_{s,t} = 0\) （若 \(t > s\) ）；

- 我们有 \(A_{s,t} = 0\) （若 \(1 \leqslant t < s \leqslant q\) ）；

- 对每个整数 \(s\) （满足 \(q \leqslant s \leqslant p\) ），存在 \(t\) 使得 \(A_{s,t} \neq 0\) 。

\begin{enumerate}
\def\labelenumi{\alph{enumi})}
\setcounter{enumi}{1}
\item
  确定这种划分的所有可能情形。（特别地，证明整数 q 与 p 都是唯一确定的，集组 \(\{I_{1},\ldots,I_{g}\}\) 与集组 \(\{I_{q+1},\ldots,I_{p}\}\) 亦然。）
\item
  存在 A 的一个特征向量 \(z \succ 0\) 属于特征值 \(\varrho\) ，当且仅当对每个整数 \(s \in [1, p]\) ，\(A_{s,s}\) 的谱半径等于 \(\varrho\) （若 \(s \leqslant q\) ），否则严格小于 \(\varrho\) 。
\item
  为使 \(A\) 与 \({}^t A\) 二者都承认一个严格正的极端向量，必须且只需存在如上的一个划分，使得 \(A_{s,t} = 0\) （对 \(1 \leqslant t < s\) 且 \(q < s \leqslant p\) ），且矩阵 \(A_{s,s}\) （\(s \leqslant q\) ）的谱半径等于 \(\varrho\) 。
\item
  若 A 与 \(^{t}A\) 都承认一个严格正的极端向量，且 \(\varrho\) 的谱重数等于 1，则 A 是不可约的。
\end{enumerate}

\begin{enumerate}
\def\labelenumi{\arabic{enumi})}
\setcounter{enumi}{21}
\tightlist
\item
  设 A 是 (I, I) 型的正且不可约的实方阵；\(\varrho\) 是它的谱半径。设 \(d \geqslant 1\) 是由习题 14 确定的整数。当 d = 1 时，我们称 A 是本原的。设 \(\Delta(\mathrm{X})\) 是多项式 \(\det(\mathrm{X} \cdot 1_{\mathrm{I}} - A)\) 。
\end{enumerate}

\begin{enumerate}
\def\labelenumi{\alph{enumi})}
\item
  我们写 \(\Delta(\mathrm{X}) = \mathrm{X}^{n} + a_{1}\mathrm{X}^{n_{1}} + \cdots + a_{t}\mathrm{X}^{n_{t}}\) ，其中 \(a_{1}, \ldots, a_{t}\) 是非零实数，\(n_{1}, \ldots, n_{t}\) 是满足 \(n > n_{1} > \cdots > n_{t}\) 的自然数。证明 d 是 \(n - n_{1}, n_{1} - n_{2}, \ldots, n_{t-1} - n_{t}\) 的最大公因子。
\item
  存在非零整数 m 与多项式 g，使得 \(\Delta(\mathrm{X}) = g(\mathrm{X}^{d})\mathrm{X}^{m}\) 。
\item
  若存在整数 p \textgreater{} 0 使得 \(A^{p}\) 严格正，则 A 是本原的。
\item
  假设 \(A\) 是本原的。引入习题 15 的矩阵 \(C(\mathbf{X})\) ，并按该习题的记号置 \(\mathsf{X}(\mathbf{X}) = \Delta (\mathbf{X}) / \Delta_1(\mathbf{X})\) 。序列 \((A^k /\varrho^k)_{k\in \mathbf{N}}\) 的极限等于 \(C(\varrho) / \mathsf{X}'(\varrho)\) 。（可利用 \(A\) 的若尔当块分解。）利用习题 15 的结果，推出存在整数 \(p > 0\) 使得 \(A^p\) 严格正。
\item
  假设存在整数 q \textgreater{} 0 使得 \(A^{q}\) 是可约的。设 p 是 q 与 d 的最大公因子。证明：存在 I 的一个分成非空部分的划分 \((\mathrm{I}_{s})_{1 \leqslant s \leqslant p}\) ，使 \(A^{q}\) 关于此划分是分块对角的，其对角块不可约且谱半径为 \(\varrho^{q}\) 。
\item
  若 A 是本原的，则 \(A^{p}\) 对每个 p \textgreater{} 0 都是不可约的。
\end{enumerate}

\(g)\) 存在 I 的一个分成非空部分的划分 \((\mathrm{I}_s)_{1 \leqslant s \leqslant d}\) ，使 \(A^d\) 关于此划分是分块对角的，其对角块是本原的且谱半径等于 \(\varrho^d\) 。

\begin{enumerate}
\def\labelenumi{\arabic{enumi})}
\setcounter{enumi}{22}
\tightlist
\item
  设 P 是 (I, I) 型的实方阵；若 P 是非负的，且向量 \((1, \ldots, 1) \in \mathbf{R}^{\mathrm{I}}\) 是 P 的属于特征值 1 的特征向量，则称 P 是随机的。
\end{enumerate}

设 A 是 (I, I) 型的实方阵。证明：A 是正的且承认一个严格正的极端向量，当且仅当存在一个随机矩阵 P、一个实数 r \textgreater{} 0 以及一个对角系数严格为正的对角矩阵 Z，使得 \(A = rZPZ^{-1}\) 。在这种情形，\(A\) 的谱半径等于 \(r\) ，而 \(A\) 的一个极端向量是 \(Z(1, \ldots, 1)\) 。

\begin{enumerate}
\def\labelenumi{\arabic{enumi})}
\setcounter{enumi}{23}
\tightlist
\item
  \begin{enumerate}
  \def\labelenumii{\alph{enumii})}
  \tightlist
  \item
    设 P 是 n 阶的实随机方阵。实数 1 是 P 的极小多项式的单根。（可利用习题 20。）此外，P 的极小多项式的每个模为 1 的根都是单根。
  \end{enumerate}
\end{enumerate}

\begin{enumerate}
\def\labelenumi{\alph{enumi})}
\setcounter{enumi}{1}
\tightlist
\item
  设 A 是一个正实矩阵，它承认一个极端向量 \(z \succ 0\) ，谱半径为 \(\varrho\) 。A 的极小多项式的所有模为 \(\varrho\) 的根都是单根。
\end{enumerate}

在习题 25 至 30 中，我们考虑一个实弗雷歇空间 E 和 E 中的一个闭凸锥 C，并假设 C 的内部非空。由 C 诱导的 E 上的序关系记作 \(\preccurlyeq_{\mathrm{C}}\) 。于是，对 E 中的 \(x\) 、\(y\) ，\(x \succcurlyeq_{\mathrm{C}} y\) 当且仅当 \(x - y \in \mathrm{C}\) 。我们将用 \(x \succ_{\mathrm{C}} y\) 表示关系 \(x - y \in \mathring{\mathrm{C}}\) 。

我们考虑满足下列条件的连续映射 u: E → E：

\begin{enumerate}
\def\labelenumi{(\roman{enumi})}
\item
  映射 u 是正齐次的：\(u(\lambda x) = \lambda u(x)\) 对每个实数 \(\lambda \geqslant 0\) 成立；
\item
  映射 u 关于 C 是递增的，即 \(x \succsim_{C} y\) 蕴含 \(u(x) \succsim_{C} u(y)\) ；
\item
  对每个 \(x \in \mathrm{C} - \{0\}\) ，有 \(u(x) \in \mathring{\mathrm{C}}\) 。
\end{enumerate}

这些习题取自 P.-L. 利翁斯 2020/2021 学年在法兰西公学院的课程，其目标是在某些条件下证明：存在元素 \(x \in \mathrm{C} - \{0\}\) 与 \(\lambda > 0\) 使得 \(x = \lambda u(x)\) ；特别地，这给出了克赖因–鲁特曼定理（III, p. 94, 定理 4 与 III, p. 97, 命题 8）的另一个证明。

对 \(f \in \mathring{\mathrm{C}}\) 与 \(\lambda \in \mathbf{R}_{+}\) ，我们把方程 \(y = \lambda u(y) + f\) 的上解（相应地，下解）称为一个元素 \(y \in \mathrm{C}\) ，它满足 \(y \succcurlyeq_{\mathrm{C}} \lambda u(y) + f\) （相应地，满足 \(y \preccurlyeq_{\mathrm{C}} \lambda u(y) + f\) ）。我们用 \(U_{f}\) 表示由 \(\lambda \in \mathbf{R}_{+}\) （使方程 \(y = \lambda u(y) + f\) 的上解集合非空）所成的集。

\begin{enumerate}
\def\labelenumi{\arabic{enumi})}
\setcounter{enumi}{24}
\tightlist
\item
  \begin{enumerate}
  \def\labelenumii{\alph{enumii})}
  \tightlist
  \item
    证明：对每个 \(x \in C\) 与每个 \(y \in \mathring{C}\) ，存在实数 \(\lambda \geqslant 0\) 使得 \(x \preccurlyeq_{C} \lambda y\) 。
  \end{enumerate}
\end{enumerate}

\begin{enumerate}
\def\labelenumi{\alph{enumi})}
\setcounter{enumi}{1}
\tightlist
\item
  集合 \(U_{f}\) 是包含 0 的一个区间。
\end{enumerate}

我们对映射 u 引入下列条件。

\begin{enumerate}
\def\labelenumi{(\roman{enumi})}
\setcounter{enumi}{3}
\item
  对每个 \(\lambda \in \mathbf{R}_{+}\) 、每个 \(f \in \mathring{\mathrm{C}}\) ，以及每个序列 \((x_{n})\) （由 C 中元素组成，满足 \(x_{n+1} = \lambda u(x_{n}) + f\) （对所有 \(n\) ），且递增并上有界或递减），它都是收敛的；
\item
  E 的每个有界子集在 u 下的像是相对紧的。
\item
  对每个映射 \(\lambda \mapsto \varepsilon_{\lambda}\) （从 \(U_f\) 到 C 中），若它满足 \(\lim_{\lambda \to \lambda_m} \varepsilon_{\lambda} = 0\) ，则集合
\end{enumerate}

\[
\{z \in \mathrm{C} \mid \exists \lambda \in \mathrm{U}_{f}\text{ 使得 } z = \lambda u(z) + \varepsilon_{\lambda} \}
\]

在 E 中是相对紧的。

\begin{enumerate}
\def\labelenumi{\arabic{enumi})}
\setcounter{enumi}{25}
\tightlist
\item
  我们沿用前面的记号与假设，并另外假设条件 (iv) 成立。
\end{enumerate}

\begin{enumerate}
\def\labelenumi{\alph{enumi})}
\item
  设 \(\lambda \in \mathrm{U}_f\) 。我们定义一个序列 \((x_n)_{n \in \mathbf{N}}\) ：\(x_0 = 0\) 且 \(x_{n+1} = \lambda u(x_n) + f\) （对所有 \(n \in \mathbf{N}\) ）。它递增且上有界；我们有 \(x_n \in \mathring{\mathrm{C}}\) （对 \(n \geqslant 1\) ）。
\item
  由条件 (iv)，序列 \((x_{n})\) 收敛；用 \(x(\lambda)\) 表示它的极限。证明 \(x(\lambda) \in \mathring{\mathrm{C}}\) 与 \(x(\lambda) = \lambda u(x(\lambda)) + f\) 。
\item
  设 \(\lambda \in \mathrm{U}_f\) ，并设 \(y \in \mathrm{C}\) 满足 \(y \succcurlyeq_{\mathrm{C}} \lambda u(y) + f\) ；由 \(x_0 = y\) 、\(x_{n+1} = \lambda u(x_n) + f\) 所定义的序列是递减的，取值于 C 中，其极限 \(x\) 满足方程 \(x = \lambda u(x) + f\) 。
\end{enumerate}

\(d)\) \(\mathrm{U}_{f}\) 在 \(\mathbf{R}_{+}\) 中的上确界不属于 \(\mathrm{U}_{f}\) 。

\begin{enumerate}
\def\labelenumi{\alph{enumi})}
\setcounter{enumi}{4}
\item
  \(U_{g} = U_{f}\) 对所有 \(g \in \mathring{C}\) 成立（对每个 \(\lambda \in U_{f}\) ，设 \(x \in C\) 满足 \(x \succsim_{C} \lambda u(x) + f\) ；则存在 M \textgreater{} 0 使得 y = Mx 满足 \(y \succsim_{C} \lambda u(y) + g\) 。） f) 设 \(\lambda \leqslant \mu\) 是实数。设 \(x\) 是方程 \(y = \lambda u(y) + f\) 的一个下解，设 \(x'\) 是方程 \(y = \mu u(y) + f\) 的一个上解。则 \(x \preccurlyeq_{C} x'\) 。（证明由 \(\varepsilon \in [0,1]\) （使 \(\varepsilon x \preccurlyeq_{C} x'\) ）组成的集与 {[}0,1{]} 重合。）
\item
  若 C 是尖锥，则对每个 \(\lambda \in \mathrm{U}_f\) ，存在唯一的元素 \(x \in \mathring{\mathrm{C}}\) 使得 \(x = \lambda u(x) + f\) 。
\end{enumerate}

\(h\) ) 设 \(\mu\in\mathbf{R}_{+}\) 满足 \(f\preccurlyeq_{\mathrm{C}}\mu u(f)\) 。则 \(\mu\notin\mathrm{U}_{f}\) 。（序列 \((x_{n})_{n\in\mathbf{N}}\) （由 \(x_{0}=0\) 与 \(x_{n+1}=\mu u(x_{n})+f\) 定义）满足 \(x_{n+1}-x_{n}\succcurlyeq_{\mathrm{C}}f\) （对所有 \(n\in\mathbf{N}\) ）。）特别地，区间 \(\mathrm{U}_{f}\) 有界。此后我们用 \(\lambda_{\mathrm{m}}\) 表示 \(\mathrm{U}_{f}\) 的上确界。

\begin{enumerate}
\def\labelenumi{\roman{enumi})}
\tightlist
\item
  设 \(z \in \mathrm{C} - \{0\}\) 与 \(\mu \in \mathbf{R}_+\) 满足 \(\mu u(z) \succsim_{\mathrm{C}} z\) 。则 \(\lambda_{\mathrm{m}} \leqslant \mu\) ，从而
\end{enumerate}

\[
\lambda_ {\mathrm{m}} \leqslant \inf _ {z \in \mathrm{C} - \{0 \}} \inf \{\mu \mid z \leqslant \mu u (z) \}.
\]

\begin{enumerate}
\def\labelenumi{\alph{enumi})}
\setcounter{enumi}{9}
\tightlist
\item
  映射 \(\lambda \mapsto x(\lambda)\) （其中 \(x(\lambda)\) 在问题 \(a\) 中定义）是递增的。
\end{enumerate}

\(k)\) 映射 \(\lambda \mapsto x(\lambda)\) 当 \(\lambda\) 趋于 \(\lambda_{\mathrm{m}}\) 时没有极限。更确切地说，若 \((\lambda_n)_n\) 是 \(\mathrm{U}_f\) 中趋于 \(\lambda_{\mathrm{m}}\) 的一个序列，则序列 \((x(\lambda_n))_n\) 不收敛。

\begin{enumerate}
\def\labelenumi{\arabic{enumi})}
\setcounter{enumi}{26}
\tightlist
\item
  我们假设 (i) 至 (v) 各条假设成立。
\end{enumerate}

\begin{enumerate}
\def\labelenumi{\alph{enumi})}
\item
  证明：\(x(\lambda)\) （对 \(\lambda \in \mathrm{U}_f\) ）所成的集不是有界的。
\item
  设 d 是 E 上定义其拓扑的一个距离。对 \(\lambda \in R_{+}\) ，定义 \(\hat{x}(\lambda) = d(0, x(\lambda))^{-1}x(\lambda)\) 。证明映射 \(\hat{x}\) 有界，且 \(\hat{x}\) 的值集有一个聚点 \(x_{m}\) 满足
\end{enumerate}

\begin{enumerate}
\def\labelenumi{(\roman{enumi})}
\item
  \(d(0, x_{\mathrm{m}}) = 1,\)
\item
  \(x_{\mathrm{m}} = \lambda_{\mathrm{m}} u(x_{\mathrm{m}}),\)
\item
  \(x_{m} \in \overset{\circ}{C}.\)
\end{enumerate}

\begin{enumerate}
\def\labelenumi{\arabic{enumi})}
\setcounter{enumi}{27}
\tightlist
\item
  在本习题中我们假设条件 (i) 至 (iv) 成立。 a) 我们固定 \(x_0 \in \mathring{\mathrm{C}}\) 。对 \(\lambda \in \mathrm{U}_f\) ，令 \(\mathrm{A}(\lambda) = \inf \{\mathrm{M} \in \mathbf{R}_+ \mid x(\lambda) \preccurlyeq_{\mathrm{C}} \mathrm{M}x_0\}\) 。设 \(\mathrm{C}_0 \in \mathbf{R}_+\) 满足 \(f \preccurlyeq_{\mathrm{C}} \mathrm{C}_0 u(x_0)\) 。对每个 \(\lambda\) （\(\mathrm{U}_f\) 中的），区间 \([\lambda, \lambda + \mathrm{C}_0 / \mathrm{A}(\lambda)[\) 包含于 \(\mathrm{U}_f\) （参见习题 26 的问题 \(d\) )）。
\end{enumerate}

\begin{enumerate}
\def\labelenumi{\alph{enumi})}
\setcounter{enumi}{1}
\tightlist
\item
  由此推出：函数 \(\lambda \mapsto \mathrm{A}(\lambda)\) 趋于 \(+\infty\) （当 \(\lambda \in \mathrm{U}_f\) 趋于 \(\lambda_{\mathrm{m}}\) 时）。
\end{enumerate}

对每个 \(\lambda\) （\(\mathrm{U}_f\) 中的），定义 \(\hat{x} (\lambda) = x(\lambda) / \mathrm{A}(\lambda)\) 。

我们现在假设条件 (vi) 成立。

\begin{enumerate}
\def\labelenumi{\alph{enumi})}
\setcounter{enumi}{2}
\tightlist
\item
  证明映射 \(\hat{x}\) 的值集有一个聚点 \(x_{\mathrm{m}}\) 满足
\end{enumerate}

\[
0 \preccurlyeq_ {\mathrm{C}} x _ {\mathrm{m}} \preccurlyeq_ {\mathrm{C}} x _ {0} \quad \text {and} \quad x _ {\mathrm{m}} = \lambda_ {\mathrm{m}} u (x _ {\mathrm{m}}).
\]

\begin{enumerate}
\def\labelenumi{\alph{enumi})}
\setcounter{enumi}{3}
\item
  证明 \(x_{m} \neq 0\) 。（若 \(x_{m} = 0\) ，设 U 是 E 中 0 的一个邻域，使得 \(x_{0} + U \subset C\) ；则对每个 \(\mu \geqslant 0\) ，将存在一个序列 \((\lambda_{n})\) （\(U_{f}\) 中的），收敛于 \(\lambda_{m}\) ，使 \(x_{0} - \mu \hat{x}(\lambda_{n}) \succsim_{C} 0\) ，这与 \(A(\lambda)\) 的定义矛盾。）
\item
  设 \(y \in \mathrm{C} - \{0\}\) 满足 \(y = \lambda_{\mathrm{m}} u(y)\) 。置 \(\theta = \sup \{\gamma \mid \gamma y \preccurlyeq_{\mathrm{C}} x_{\mathrm{m}}\}\) 。我们有 \(\theta \in \mathbf{R}_{+}^{*}\) 且 \(x_{\mathrm{m}} = \theta y\) 。（我们有 \(\theta y \preccurlyeq_{\mathrm{C}} x_{\mathrm{m}}\) 与 \(\theta y = \lambda_{\mathrm{m}} u(\theta y)\) ；若 \(x_{\mathrm{m}} \neq \theta y\) ，则将有 \(x_{\mathrm{m}} - \theta y \succ_{\mathrm{C}} 0\) ，且存在 \(\varepsilon > 0\) 使 \(x_{\mathrm{m}} - \theta y \succcurlyeq_{\mathrm{C}} \varepsilon y\) ，这与 \(\theta\) 的定义矛盾。）
\item
  设 \(y \in C - \{0\}\) 与 \(\mu \in R\) 满足 \(y = \mu u(y)\) 。则 \(\mu = \lambda_{m}\) 。（注意 \(0 < \mu\) 且 \(\mu \geqslant \lambda_{m}\) ；若有 \(\mu > \lambda_{m}\) ，令 \(\theta = \sup\{\gamma \mid \gamma x_{m} \preccurlyeq_{C} y\}\) ；则将有 \(\theta \in R_{+}^{*}\) ，进而 \(\theta x_{m} \preccurlyeq_{C} y\) ，于是 \(\theta x_{m} = \lambda_{m} u(\theta x_{m}) \prec_{C} \mu u(\theta x_{m}) \preccurlyeq_{C} \mu u(y) = y\) ，这与 \(\theta\) 的定义矛盾。）
\end{enumerate}

\begin{enumerate}
\def\labelenumi{\arabic{enumi})}
\setcounter{enumi}{28}
\tightlist
\item
  我们沿用前一习题的记号；特别地，假设 C 的内部非空。我们假设 u 是线性映射，且条件 (i)、(ii)、(iii) 成立。
\end{enumerate}

\begin{enumerate}
\def\labelenumi{\alph{enumi})}
\tightlist
\item
  若 u 是 E 的紧自同态，则条件 (iv) 与 (v) 成立。
\end{enumerate}

我们假设条件 (iv) 与 (vi) 成立。设 \((\lambda_{\mathrm{m}}, x_{\mathrm{m}})\) 是由前一习题给出的。

\begin{enumerate}
\def\labelenumi{\alph{enumi})}
\setcounter{enumi}{1}
\tightlist
\item
  对每个 \(\lambda \in \mathbf{R}_{+}\) （满足 \(\lambda < \lambda_{\mathrm{m}}\) ），自同态 \(1_{\mathrm{E}} - \lambda u\) 是可逆的。（首先证明 \(1_{\mathrm{E}} - \lambda u\) 通过过渡到子空间在 \(\mathring{\mathrm{C}}\) 上诱导一个双射，参见前一习题的问题 \(a\) ) 与 \(f\) )；然后推出
\end{enumerate}

\(1_{E} - \lambda u\) 是满射，然后通过把 \(x_{1} - x_{2}\) 写成 \(1_{E} - \lambda u\) 的核中元素（其中 \(x_{i} \in \mathring{C}\) ）并验证存在 \(z \in E\) 使 \(x_{1} + z - \lambda u(x_{1} + z) \in \mathring{C}\) ，推出它是单射。

\begin{enumerate}
\def\labelenumi{\alph{enumi})}
\setcounter{enumi}{2}
\tightlist
\item
  设 E 中的 y 满足 \(y = \lambda_{\mathrm{m}} u(y)\) 。则 \(y \in C\) 。（存在 c 使 \(cx_{m} + y \succ_{C} 0\) ；利用前一习题的 (e) 部分得出结论。）
\end{enumerate}

\(d)\) 设 \(\mathrm{C}^{\circ}\subset\mathrm{E}^{\prime}\) 是 C 的极集。对每个 \(x\in\mathring{\mathrm{C}}\) 与每个 \(\xi\in\mathrm{C}^{\circ}-\{0\}\) ，有 \(\langle x,\xi\rangle>0\) 。

\begin{enumerate}
\def\labelenumi{\alph{enumi})}
\setcounter{enumi}{4}
\tightlist
\item
  C 的极集 C° 在 \(^{t}u\colon E'\to E'\) 下不变；对每个 \(\lambda\in R_{+}\) （满足 \(\lambda<\lambda_{m}\) ），自同态 \(1_{E'}-\lambda^{t}u\) （\(E'\) 的）是一个同构。
\end{enumerate}

\begin{enumerate}
\def\labelenumi{\arabic{enumi})}
\setcounter{enumi}{29}
\tightlist
\item
  我们沿用前一习题的假设，即 u 是线性、紧的，且假设条件 (i) 至 (v) 成立。自同态 \(^{t}u\) 也是紧的。
\end{enumerate}

\begin{enumerate}
\def\labelenumi{\alph{enumi})}
\item
  设 \(\eta \in \mathrm{C}^{\circ} - \{0\}\) 。对每个 \(\lambda \in \mathbf{R}_{+}\) （满足 \(\lambda < \lambda_{\mathrm{m}}\) ），唯一解 \(\xi(\lambda)\) （方程 \(\xi - \lambda^{t}u(\xi) = \eta\) 的）属于 \(\mathrm{C}^{\circ}\) 。（若 \(x \in \mathrm{C}\) ，设 \(y \in \mathrm{C}\) 满足 \(y - \lambda u(y) = x\) ；则 \(\langle y, \eta \rangle = \langle x, \xi \rangle\) 。）
\item
  设 \(x_{0} \in \mathring{C}\) 。映射 \(\lambda \mapsto \langle x_{0}, \xi(\lambda) \rangle\) 当 \(\lambda\) 趋于 \(\lambda_{m}\) 时趋于无穷。证明值 \(\xi(\lambda)/\langle x_{0}, \xi(\lambda) \rangle\) 所成的集有一个聚点 \(\xi_{m}\) 满足：
\end{enumerate}

\begin{enumerate}
\def\labelenumi{(\roman{enumi})}
\item
  \(\xi_{\mathrm{m}}\in \mathrm{C}^{\circ}\) ;
\item
  \(\langle x_{0},\xi_{m}\rangle=1;\)
\item
  \(\xi_{\mathrm{m}} = \lambda_{\mathrm{m}}^{t}u(\xi_{\mathrm{m}})\) .
\end{enumerate}

\begin{enumerate}
\def\labelenumi{\alph{enumi})}
\setcounter{enumi}{2}
\item
  设 \(f \in E\) 满足 \(\langle f, \xi_{m} \rangle = 0\) 。对 \(\lambda \in R_{+}\) （满足 \(\lambda < \lambda_{m}\) ），用 \(x(\lambda)\) 表示唯一解 \(x\) （方程 \(x - \lambda u(x) = f\) 在 E 中的）。证明 \(x(\lambda)\) 所成的集在 E 中有界（设 d 是定义 E 的拓扑的距离；若存在序列 \(\lambda_{n} \to \lambda_{m}\) 使 \(x(\lambda_{n})\) 趋于无穷，置 \(\hat{x}_{n} = x(\lambda_{n}) / d(0, x(\lambda_{n}))\) ，并证明序列 \((\hat{x}_{n})\) 收敛于一个元素 \(\hat{x}\) ，它满足 \(d(0, \hat{x}) = 1\) 、\(\hat{x} = \lambda u(\hat{x})\) 与 \(\langle \hat{x}, \xi_{m} \rangle = 0\) ，这与必有 \(\hat{x} \in \mathring{C}\) 的事实矛盾）。
\item
  存在唯一的元素 \(x \in E\) 满足 \(f = x - \lambda_{\mathrm{m}} u(x)\) 与 \(\langle x, \xi_{\mathrm{m}} \rangle = 0\) 。（关于存在性，利用前一问题证明存在一个序列 \((\lambda_{n})\) 收敛于 \(\lambda_{m}\) ，使 \((x(\lambda_{n}))\) 收敛于一个满足所需条件的元素 x。）
\item
  自同态 \(1_{\mathrm{E}'} - \lambda_{\mathrm{m}}{}^{t}u\) （\(\mathrm{E}'\) 的）通过过渡到子空间诱导极集 \(x_{\mathrm{m}}^{\circ}\) （\(x_{\mathrm{m}}\) 的）上的一个双射。
\item
  特征空间 \(\operatorname{Ker}(1_{\mathrm{E}} - \lambda_{\mathrm{m}}u)\) 的维数为 1，且等于 \(1_{E} - \lambda_{m}u\) 的像的极集。
\end{enumerate}

第四章

\chapter{希尔伯特谱理论}
\setcounter{exercisepar}{0}

对每个向量空间 E，我们用 \(1_{E}\) 表示 E 的恒等映射。若 E、F 与 G 是向量空间，且 v: \(E \rightarrow F\) 与 u: \(F \rightarrow G\) 是线性映射，则从 E 到 G 的线性映射 \(u \circ v\) 有时记作 uv。

设 E 是一个预希尔伯特空间。我们用 \(\langle x|y\rangle_{E}\) ，或简作 \(\langle x|y\rangle\) ，表示 E 上的内积。E 的子集 A 的正交记作 A°。它是 E 的一个闭子空间；它化为 0 当且仅当 A 在 E 中是全的（EVT, V, p. 16, 推论 1）。若 F ⊂ E 是 E 的向量子空间，则 F°° = \(\overline{F}\) （参见 EVT, V, p. 13）。E 的自同态的谱总是相对于巴拿赫代数 \(\mathcal{L}(E)\) 而言的。

设 F 是希尔伯特空间。我们用 \(\mathcal{L}_1(\mathrm{E})\) 表示 E 的有限迹自同态所成的空间（EVT, V, p. 50, 定义 8），用 \(\mathcal{L}_2(\mathrm{E};\mathrm{F})\) 表示从 E 到 F 的希尔伯特–施密特映射所成的空间（EVT, V, p. 51, 定义 9）。映射 \(u\mapsto \| u\| _2 = (\mathrm{Tr}(u^{*}u))^{1 / 2}\) 是 \(\mathcal{L}_2(\mathrm{E};\mathrm{F})\) 上的一个希尔伯特范数（EVT, V, p. 52, 定理 1, (i)）。

有限秩的连续线性映射所成的空间 \(\mathcal{L}^{\mathrm{f}}(\mathrm{E};\mathrm{F})\) 在 \(\mathcal{L}^{\mathrm{c}}(\mathrm{E};\mathrm{F})\) 中稠密（III, p. 15, 命题 10 与 III, p. 16, 命题 11 的推论）。特别地，空间 \(\mathcal{L}_2(\mathrm{E};\mathrm{F})\) 在 \(\mathcal{L}^{\mathrm{c}}(\mathrm{E};\mathrm{F})\) 中稠密。

E 与 F 的希尔伯特张量积（EVT, V, p. 28, 定义 1）记作 E \(\widehat{\otimes}_2\) F。

若 X 是局部紧拓扑空间，\(\mu\) 是 X 上的一个测度，我们记 \(\mathcal{L}^{p}(\mathrm{X},\mu)=\mathcal{L}_{\mathbf{C}}^{p}(\mathrm{X},\mu)\) 与 \(\mathrm{L}^{p}(\mathrm{X},\mu)=\mathrm{L}_{\mathbf{C}}^{p}(\mathrm{X},\mu)\) （对 \(p\in[1,+\infty]\) ）。

\section{希尔伯特空间上的紧算子}

在本节中，K 是等于 R 或 C 的域，E 表示域 K 上的一个希尔伯特空间。

\subsection{对角自同态}

定义 1. --- 设 \(\mathrm{B} = (e_i)_{i \in \mathrm{I}}\) 是 E 的一个标准正交基。称一个自同态 \(u\) 在基 B 中是对角的或相对于 B 是对角的，如果存在由 K 的元素组成的族 \(\lambda = (\lambda_i)_{i \in \mathrm{I}}\) ，使得 \(u(e_i) = \lambda_i e_i\) （对所有 \(i \in \mathrm{I}\) ）。

设 \(\mathrm{B} = (e_i)_{i \in \mathrm{I}}\) 是 E 的一个标准正交基。我们用 \(\mathcal{D}_{\mathrm{B}}(\mathrm{E})\) 表示 E 的相对于 B 为对角的自同态所成的集。这是 \(\mathcal{L}(\mathrm{E})\) 的一个闭的、幺的交换子代数。设 \(u \in \mathcal{D}_{\mathrm{B}}(\mathrm{E})\) 。族 \(\lambda = (\lambda_i)_{i \in \mathrm{I}}\) （满足 \(u(e_i) = \lambda_i e_i\) ）由 u 和基 B 唯一确定，称为 u 相对于 B 的特征值族。

假设 K = R，并把 E 与 \(\mathrm{E}_{(\mathbf{C})}\) 的一个子空间等同。设 \(\mathrm{B} = (e_i)_{i \in \mathrm{I}}\) 是 E 的一个标准正交基。用 \(\mathrm{B}_{(\mathbf{C})}\) 表示标准正交基 \((1 \otimes e_i)_{i \in \mathrm{I}}\) （\(\mathrm{E}_{(\mathbf{C})}\) 的）（参见 EVT, V, p. 29, 推论 1）。映射 \(u \mapsto u_{(\mathbf{C})}\) 诱导从 \(\mathcal{D}_{\mathrm{B}}(\mathrm{E})\) 到 \(\mathcal{D}_{\mathrm{B}_{(\mathbf{C})}}(\mathrm{E}_{(\mathbf{C})})\) 中的一个单的代数态射；它的像是 \(u \in \mathcal{D}_{\mathrm{B}_{(\mathbf{C})}}(\mathrm{E}_{(\mathbf{C})})\) （满足 \(u(\mathrm{E}) \subset \mathrm{E}\) ）所成的集；换言之，即基 \(\mathrm{B}_{(\mathbf{C})}\) 中特征值为实的对角自同态 u 所成的集。

对每个 \(i \in \mathrm{I}\) ，用 \(p_i\) 表示 E 的以 \(\mathrm{Ke}_i\) 为像的正交投影。我们有 \(\| p_i \| = 1\) （对所有 \(i \in \mathrm{I}\) ）。自同态 \(p_i\) 在基 B 中是对角的，其特征值族是族 \((\delta_{ij})_{j\in\mathrm{I}}\) （克罗内克符号，参见 A, II, p. 24）。

设 \(\mathrm{B}^{\prime}=(f_{i})_{i\in\mathrm{I}}\) 是 E 的一个标准正交基，\(u:E\to E\) 是满足 \(u(f_{i})=e_{i}\) 的等距同构。则 \(\mathcal{D}_{\mathrm{B}^{\prime}}(\mathrm{E})=u^{-1}\mathcal{D}_{\mathrm{B}}(\mathrm{E})u\) 。

我们给 I 赋以离散拓扑，并用 \(\mathcal{B}(\mathrm{I}) = \mathcal{C}_{b}(\mathrm{I}; \mathrm{K})\) 表示 I 上取值于 K 的有界函数所成的幺巴拿赫代数（I, p. 17, 例 2）；若 K = C，它是一个星代数（I, p. 102, 例 2）。

引理 1. --- 设 \(u \in \mathcal{D}_{\mathrm{B}}(\mathrm{E})\) ，\(\lambda = (\lambda_i)_{i \in \mathrm{I}}\) 是它的特征值族。族 \(\lambda\) 有界，且有 \(\sup_i |\lambda_i| = \| u\|\) 。

有 \(|\lambda_i| \leqslant \|u\|\) （对所有 \(i\) ），从而 \(\sup_{i \in I} |\lambda_i| \leqslant \|u\|\) 。此外，由于

\[
\| u (x) \| ^ {2} = \sum_ {i \in \mathrm{I}} | \lambda_ {i} | ^ {2} | \langle e _ {i} | x \rangle | ^ {2} \leqslant \left(\sup _ {i \in \mathrm{I}} | \lambda_ {i} | ^ {2}\right) \| x \| ^ {2}
\]

对所有 \(x \in \mathrm{E}\) 成立（EVT, V, p. 22, 命题 5），由此得出 \(\| u\| = \sup |\lambda_i|\) 。

命题 1. --- a) 映射 \(\alpha\) （从 \(\mathcal{D}_{\mathrm{B}}(\mathrm{E})\) 到 \(\mathcal{B}(\mathrm{I})\) 中的）把对角自同态 \(u\) 对应到 \(u\) 的特征值族，它是 K 上的巴拿赫代数的一个等距同构。若 K = C，它是对合代数的一个态射。

\begin{enumerate}
\def\labelenumi{\alph{enumi})}
\setcounter{enumi}{1}
\item
  对任何有界族 \(\lambda = (\lambda_i)_{i\in \mathrm{I}}\) ，族 \((\lambda_ip_i)_{i\in \mathrm{I}}\) 在赋以逐点收敛拓扑的空间 \(\mathcal{L}(\mathrm{E})\) 中是可和的，且该族之和是唯一的映射 \(u\in \mathcal{D}_{\mathrm{B}}(\mathrm{E})\) ，它满足 \(\alpha (u) = \lambda\) ；
\item
  设 \(u \in \mathcal{D}_{\mathrm{B}}(\mathrm{E})\) ，\(\lambda\) 是它相对于 B 的特征值族。\(u\) 的谱是 \(\lambda\) 的值集在 K 中的闭包；
\item
  若 K = C，则从 \(\mathcal{C}(\mathrm{Sp}(u))\) 到 \(\mathcal{L}(\mathrm{E})\) 中的连续函数演算映射把连续函数 \(f \in \mathcal{C}(\mathrm{Sp}(u))\) 对应到自同态 \(\alpha(f \circ \lambda)\) （\(\mathcal{D}_{\mathrm{B}}(\mathrm{E})\) 中的）。
\end{enumerate}

由引理 1，在 B 中为对角的自同态的特征值族是有界的，且如此定义的从 \(\mathcal{D}_{\mathrm{B}}(\mathrm{E})\) 到 \(\mathcal{B}(\mathrm{I})\) 中的映射是幺巴拿赫代数的一个连续等距态射。

设 \(i \in I\) 。对每个 \(j \in I\) ，我们有 \(\langle u^{*}(e_{i}) | e_{j} \rangle = \lambda_{j} \langle e_{i} | e_{j} \rangle = \langle \overline{\lambda_{j}} e_{i} | e_{j} \rangle\) 。由此得出 \(u^{*}(e_{i}) = \overline{\lambda}_{i} e_{i}\) 。于是 u 的伴随是基 B 中的对角自同态，其特征值族是 \((\overline{\lambda}_{i})_{i \in I}\) 。断言 a) 由此得出。

设 \(\lambda = (\lambda_i)_{i\in \mathrm{I}}\in \mathcal{B}(\mathrm{I})\) 。对每个 \(x\in \mathrm{E}\) ，族 \((|\langle e_i|x\rangle |^2)_{i\in \mathrm{I}}\) 可和，其和为 \(\| x\|^2\) （EVT, V, p. 22, 命题 5），从而，对 I 的每个有限子集 J：

\[
\begin{array}{r l r} & & {\left\| \sum_ {i \in \mathrm{J}} \lambda_ {i} p _ {i} (x) \right\| ^ {2} = \sum_ {i \in \mathrm{J}} | \lambda_ {i} | ^ {2} | \langle e _ {i} | x \rangle | ^ {2} \leqslant \Bigl (\underset {i \in \mathrm{I}} {\sup} | \lambda_ {i} | ^ {2} \Bigr) \sum_ {i \in \mathrm{J}} | \langle e _ {i} | x \rangle | ^ {2}} \\ & & {\leqslant \Bigl (\underset {i \in \mathrm{I}} {\sup} | \lambda_ {i} | \Bigr) ^ {2} \| x \| ^ {2}.} \end{array}
\]

因此，族 \((\lambda_{i}p_{i})\) 在赋以简单收敛拓扑的 \(\mathcal{L}(\mathrm{E})\) 中可和。它的和 \(u_{\lambda}\) 满足 \(u_{\lambda}(e_{i}) = \lambda_{i}e_{i}\) ；于是它是基 B 中的对角自同态，以 \(\lambda\) 为特征值族。断言 b) 由此得出。

最后的断言由 \(a\) )、I, p. 17 的例 3 与 I, p. 111 的例 4 得出。

注记. --- 巴拿赫代数 \(\mathcal{D}_{\mathrm{B}}(\mathrm{E})\) 是 \(\mathcal{L}(\mathrm{E})\) 的一个极大交换闭子代数。事实上，设 u 是与 \(\mathcal{D}_{\mathrm{B}}(\mathrm{E})\) 交换的 E 的自同态。设 \(i \in I\) 。由于正交投影 \(p_{i}\) 在基 B 中是对角的，我们有 \(p_{i}(u(e_{i})) = u(p_{i}(e_{i})) = u(e_{i})\) ，这蕴含 \(u(e_{i})\) 与 \(e_{i}\) 成比例。于是 u 在基 B 中是对角的。

若 E 是无穷维的，则存在 \(\mathcal{L}(\mathrm{E})\) 的极大交换对合子代数，它们不同构于 \(\mathcal{D}_{\mathrm{B}}(\mathrm{E})\) （IV, p. 314, 习题 5）。

命题 2. --- 设 \(u \in \mathcal{D}_{\mathrm{B}}(\mathrm{E})\) ，\(\lambda = (\lambda_i)_{i \in \mathrm{I}}\) 是它的特征值族。

\begin{enumerate}
\def\labelenumi{\alph{enumi})}
\tightlist
\item
  下列条件等价：
\end{enumerate}

\begin{enumerate}
\def\labelenumi{(\roman{enumi})}
\item
  有 \(\lambda \in \mathcal{C}_{0}(\mathrm{I};\mathrm{K})\) ；
\item
  自同态 u 是紧的；
\item
  族 \((\lambda_{i}p_{i})_{i\in\mathrm{I}}\) 在巴拿赫空间 \(\mathcal{L}(\mathrm{E})\) 中可和。其和这时等于 u。
\end{enumerate}

\begin{enumerate}
\def\labelenumi{\alph{enumi})}
\setcounter{enumi}{1}
\tightlist
\item
  假设 \(u\) 是紧的，用 \(\Lambda\) 表示 \(\lambda\) 的值集。由 \(i \in I\) （使 \(\lambda_i \neq 0\) ）组成的集是可数的。若 E 是无穷维的，我们有 \(\mathrm{Sp}_{\mathrm{s}}(u) = \Lambda - \{0\}\) 与 \(\mathrm{Sp}(u) = \Lambda \cup \{0\}\) 。若 E 是有限维的，则 \(\mathrm{Sp}_{\mathrm{s}}(u) = \mathrm{Sp}(u) = \Lambda\) 。
\end{enumerate}

由引理 1，对 I 的每个有限子集 J，我们有 \(\| \sum_{j\in J}\lambda_jp_j\| = \sup_{j\in J}|\lambda_j|\) 。由此，按柯西判别准则，族 \((\lambda_ip_i)\)

在 \(\mathcal{L}(\mathrm{E})\) 中可和当且仅当族 \(\lambda\) 在无穷远处趋于 0，这蕴含条件 (i) 与 (iii) 等价。

条件 (iii) 蕴含 \(u\) 是紧的（III, p. 4, 命题 2 的推论）。反之，假设自同态 \(u\) 是紧的。由于族 B 在 E 中有界，它在 \(u\) 下的像在 E 中相对紧（III, p. 2），从而是准紧的。设 \(\varepsilon > 0\) ，设 \(J\) 是 I 的一个有限子集，使得 B 在 \(u\) 下的像包含于半径为 \(\varepsilon\) 、心为 \(u(e_j)\) （\(j \in \mathrm{J}\) ）的球之并。设 \(i \in \mathrm{I} - \mathrm{J}\) 。存在 \(j \in \mathrm{J}\) 使 \(\| u(e_i) - u(e_j) \| \leqslant \varepsilon\) ，从而

\[
\left| \lambda_ {i} \right| ^ {2} \leqslant \left| \lambda_ {i} \right| ^ {2} + \left| \lambda_ {j} \right| ^ {2} = \left\| u \left(e _ {i}\right) - u \left(e _ {j}\right) \right\| ^ {2} \leqslant \varepsilon^ {2}.
\]

因此，我们有 \(\lambda \in \mathcal{C}_0(\mathrm{I};\mathrm{K})\) ，即条件 (i)。

最后，u 的谱及其敏感谱利用 \(\lambda\) 由命题 1, c) 与 III, p. 90 命题 5 计算。

注记. --- 当 I 无限时，条件 \(\lambda\in\mathcal{C}_{0}(\mathrm{I};\mathrm{K})\) 也可表述为“族 \(\lambda\) 沿 I 的有限子集的余集的滤子趋于 0”。

\subsection{紧自同态的对角化}

定理 1. --- 假设 K = C。设 u 是 E 的紧且正规的自同态。存在 E 的一个标准正交基 B，使 u 在基 B 中是对角的。

集 \(\mathrm{Sp}_{\mathrm{s}}(u)\) 是可数的，且当 E 无限时它不含 0（III, p. 90, 命题 5, \(b\) ）。对每个元素 \(\lambda \in \mathrm{Sp}_{\mathrm{s}}(u)\) ，用 \(\mathrm{N}_{\lambda}\) 表示 \(u - \lambda 1_{\mathrm{E}}\) 的幂零空间。它是有限维的（同处），且由于 \(u\) 正规，它与 \(u\) 的相对于 \(\lambda\) 的特征空间重合（EVT, V, p. 43, 命题 8 的推论）。诸空间 \(\mathrm{N}_{\lambda}\) 两两正交（I, p. 132, 第 5 目）。

设 F 是 E 的这样的子空间：它是诸空间 \(N_{\lambda}\) （对 \(\lambda \in \mathrm{Sp}_{\mathrm{s}}(u) - \{0\}\) ）的希尔伯特和。它是可数型的空间，在 u 下稳定，因为每个子空间 \(N_{\lambda}\) 在 u 下稳定。由于 \(N_{\lambda}\) 也是 \(u^{*}\) 的相对于 \(\overline{\lambda}\) 的特征空间（EVT, V, 同处），自同态 u 通过过渡到子空间诱导一个自同态 \(\widetilde{u}\) （\(F^{\circ}\) 的）。自同态 \(\widetilde{u}\) 是紧的（III, p. 5, 命题 3）且正规的（I, p. 135, 引理 4）。按构造，\(\widetilde{u}\) 的敏感谱包含于 \(\{0\}\) （事实上，\(\widetilde{u}\) 的每个特征向量也会是 u 的特征向量，从而若对应的特征值非零，它就属于某个空间 \(N_{\lambda}\) ）。于是 \(\widetilde{u}\) 的谱半径为零，从而 \(\widetilde{u}=0\) （因为 \(\widetilde{u}\) 正规；I, p. 108, 推论 1）。因此我们有 \(F^{\circ}\subset\operatorname{Ker}(u)\) 。

对每个 \(\lambda \in \mathrm{Sp}_{\mathrm{s}}(u)\) ，设 \(B_{\lambda}\) 是 \(N_{\lambda}\) 的一个标准正交基，设 \((e_j)_{j \in J}\) 是由诸 \(B_{\lambda}\) 的并组成的族；它是 F 的一个标准正交基。设 B 是 \((e_j)_{j \in J}\) 与 \(F^{\circ}\) 的一个标准正交基之并。这是 E 的一个标准正交基，且 \(u\) 在基 B 中是对角的。

推论 1. --- 设 u 是 E 的紧埃尔米特自同态。存在 E 的一个标准正交基 B，使 u 在基 B 中是对角的且其特征值为实数。

若 K = C，这立即由定理得出。假设 K = R。空间 E 是 \(\mathrm{E}_{(\mathbf{C})}\) 上的一个 R-结构（参见 A, II, p. 119）。自同态 \(u_{(\mathbf{C})}\) （\(\mathrm{E}_{(\mathbf{C})}\) 的）是紧的（III, p. 2, 注记 4）且埃尔米特的。设 \(B = (e_j)_{j \in J}\) 是 \(\mathrm{E}_{(\mathbf{C})}\) 的一个标准正交基，使得 \(u_{(\mathbf{C})} \in \mathcal{D}_{\mathrm{B}}(\mathrm{E}_{(\mathbf{C})})\) ，设 \(\lambda\) 是 \(u_{(\mathbf{C})}\) 的特征值族（定理 1）。我们有 \(\lambda \in R^{J}\) （I, p. 106, 命题 4）；由于线性映射 \(u_{(\mathbf{C})}\) 是 R-有理的，\(u_{(\mathbf{C})}\) 的相对于 \(\lambda_j\) 的特征空间对每个 \(j \in J\) 都是 R-有理的（参见 A, V, p. 60, 命题 6）。因此存在它的一个属于 E 的基，且更有理由存在 E 中的一个标准正交基。这些基之并是 E 的一个标准正交基 \(B_R\) ，使得 \(u \in \mathcal{D}_{\mathrm{B}_R}(\mathrm{E})\) 。

推论 2. --- 设 F 是希尔伯特空间，u 是从 E 到 F 的紧连续线性映射。存在 u 的始空间的一个标准正交基 \((e_{i})_{i\in\mathrm{I}}\) （\(\operatorname{Ker}(u)^{\circ}\) 的）、F 的一个标准正交族 \((f_{i})_{i\in\mathrm{I}}\) 以及一个族 \((\alpha_{i})_{i\in\mathrm{I}}\in(\mathbf{R}_{+}^{*})^{\mathrm{I}}\) ，使得 \(u(e_{i})=\alpha_{i}f_{i}\) （对所有 \(i\in I\) ）。

设 \(v = u^{*} \circ u\) 。这是 E 的一个紧的、正的（从而埃尔米特的）自同态。由推论 1，存在 E 的一个标准正交基 \((e_{j})_{j \in \mathrm{J}}\) ，使得 \(v\) 在此基中是对角的。它的特征值族 \((\lambda_{j})_{j \in \mathrm{J}}\) 包含于 \(\mathbf{R}_{+}^{\mathrm{J}}\) 。置 \(\alpha_{j} = \sqrt{\lambda_{j}}\) （对所有 \(j \in \mathrm{J}\) ）。设 I 是由 \(j \in \mathrm{J}\) （使 \(\alpha_{j} \neq 0\) ）组成的集。由于 \(v\) 紧，这是一个可数集。族 \((e_{i})_{i \in \mathrm{I}}\) 是 \(v\) 的始空间（即始空间 \(\operatorname{Ker}(u)^{\circ}\) ，\(u\) 的）的一个标准正交基（EVT, V, p. 43, 命题 8）。置 \(f_{i} = \frac{1}{\alpha_{i}} u(e_{i})\) （对 \(i\in \mathrm{I}\) ）。对 I 中任何 \(i\) 与 \(j\) ，我们有

\[
\langle f _ {i} | f _ {j} \rangle = \frac {1}{\alpha_ {i} \alpha_ {j}} \langle u (e _ {i}) | u (e _ {j}) \rangle = \frac {1}{\alpha_ {i} \alpha_ {j}} \langle v (e _ {i}) | e _ {j} \rangle = \frac {\lambda_ {i}}{\alpha_ {i} \alpha_ {j}} \langle e _ {i} | e _ {j} \rangle ,
\]

由此得出族 \((f_{i})_{i\in\mathrm{I}}\) 在 F 中是标准正交的。由于 \(u(e_{i})=\alpha_{i}f_{i}\) （对所有 \(i\in I\) ），推论得证。

定义 2. --- 沿用推论的记号，族 \((\alpha_{i})_{i\in I}\) 是 \(u\) 的相对于始空间的标准正交基 \((e_{i})_{i\in I}\) （该始空间为 \(u\) 的始空间）的奇异值族。

注记. --- 1) 该推论推广了 EVT, V, p. 54 的定理 2，后者对应于希尔伯特–施密特映射。

\begin{enumerate}
\def\labelenumi{\arabic{enumi})}
\setcounter{enumi}{1}
\tightlist
\item
  沿用推论的记号，我们有公式
\end{enumerate}

\[
u (x) = \sum_ {i \in I} \alpha_ {i} \langle e _ {i} | x \rangle f _ {i}\tag{1}
\]

对所有 \(x \in E\) 成立。

\begin{enumerate}
\def\labelenumi{\arabic{enumi})}
\setcounter{enumi}{2}
\tightlist
\item
  设 \(u\) 是 E 的紧正自同态。设 \(B = (f_i)_{i \in I}\) 是 E 的一个标准正交基，使得 \(u\) 在基 B 中是对角的（定理 1），设 \((\lambda_i)_{i \in I}\) 是 \(u\) 在此基中的特征值族。设 J 是由 \(i \in I\) （使 \(\lambda_i > 0\) ）组成的集；族 \((e_i)_{i \in J}\) 是空间 \(\text{Ker}(u)^{\circ}\) 的一个标准正交基。对每个 \(i \in J\) ，置 \(e_i = f_i\) 与 \(\alpha_i = \lambda_i\) 。我们得到 \(u(e_i) = \alpha_if_i\) ：族 \((\alpha_i)_{i \in J}\) 是 \(u\) 的相对于基 \((e_i)_{i \in J}\) 的奇异值族。
\end{enumerate}

\subsection{特征值的递减序列}

在本节中，我们假设 K = C。

我们置 \(\overline{\mathbf{N}} = \mathbf{N} \cup \{+\infty\} \subset \overline{\mathbf{R}}\) 。在本节中，若向量空间 E 不是有限维的，就说它有维数 \(+\infty \in \overline{\mathbf{N}}\) 。设 \(\mathrm{I_E} \subset \mathbf{N}\) 是 E 的满足 \(\mathrm{F} \neq \mathrm{E}\) 的有限维子空间 F 的维数所成的集。若 E 无穷维，我们有 \(\mathrm{I_E} = \mathbf{N}\) ，否则 \(\mathrm{I_E} = \{0, \ldots, \dim(\mathrm{E}) - 1\}\) 。当不致引起混淆时，我们置 \(\mathrm{I} = \mathrm{I_E}\) 。

设 \(u\) 是 E 的紧正（特别地，埃尔米特）自同态。\(u\) 的敏感谱是正实数的一个严格递减序列 \((\nu_{k})_{0\leqslant k<\mathrm{Card}(\mathrm{Sp}_{\mathrm{s}}(u))}\) 的值集（参见

III, p. 90, 命题 5）。对每个整数 \(k\) （满足 \(0 \leqslant k < \text{Card}(\text{Sp}_s(u))\) ），用 \(n_k\) 表示 \(\nu_k\) 的谱重数，其中 \(n_k \geqslant 1\) 。设 \(\text{M} \in \overline{\mathbf{N}}\) 是谱重数 \(n_k\) 之和；它是 \(u\) 的像的维数。我们有 \(\text{M} \leqslant \text{Card}(\text{I})\) 。

对 \(0 \leqslant n < M\) ，定义 \(\lambda_n(u) = \nu_k\) ，其中 \(k \geqslant 0\) 是唯一使下式成立的整数：

\[
n _ {0} + \dots + n _ {k - 1} \leqslant n <   n _ {0} + \dots + n _ {k}.
\]

置 \(\lambda_{n}(u)=0\) （若 \(n\in I\) 满足 \(n\geqslant M\) ）。这种情形只能在 I=N 且 \(\mathrm{Sp}_{\mathrm{s}}(u)\) 有限时发生（或者等价地，E 无穷维且 u 有限秩）。

序列 \((\lambda_{n}(u))_{n\in\mathrm{I}}\) 是递减的；对每个 \(\lambda\in\mathrm{Sp}_{\mathrm{s}}(u)\) ，使 \(\lambda_{n}(u)=\lambda\) 的整数 n 的个数等于 u 的特征值 \(\lambda\) 的谱重数。

我们约定俗成地说，\((\lambda_{n}(u))_{n\in\mathrm{J}}\) 是 u 的按重数重复的特征值递减序列。

命题 3. --- 设 u 是 E 的紧正自同态。存在 E 中的一个标准正交族 \((e_{n})_{n\in\mathrm{I}}\) ，使得对每个 \(x\in E\) ，有

\[
u (x) = \sum_ {n \in \mathrm{I}} \lambda_ {n} (u) \langle e _ {n} | x \rangle e _ {n}, \quad \langle x | u (x) \rangle = \sum_ {n \in \mathrm{I}} \lambda_ {n} (u) | \langle e _ {n} | x \rangle | ^ {2}.
\]

设 \(\mathrm{B} = (f_{j})_{j \in \mathrm{J}}\) 是 E 的使 u 为对角的标准正交基（IV, p. 150, 推论 1），\((\lambda_{j})_{j \in \mathrm{J}}\) 是 u 在基 B 中的特征值族。设 J' 是由 \(j \in J\) （使 \(\lambda_{j}\) 属于 u 的敏感谱）组成的集。

对每个 \(\lambda \in \mathrm{Sp}_{\mathrm{s}}(u)\) ，在整数 \(n\) （满足 \(0 \leqslant n < \mathrm{M}\) 与 \(\lambda_{n}(u) = \lambda\) ）与 \(j \in \mathrm{J}'\) （满足 \(\lambda_{j} = \lambda\) ）之间存在一个双射，因为这两个集有相同的基数，即 \(\lambda\) 的谱重数。对每个 \(\lambda\) 选取这样的双射，就定义了双射 \(\iota\) （从满足 \(0 \leqslant n < \mathrm{M}\) 的整数所成的集到集 \(\mathrm{J}'\) 的）。序列 \((e_n)_{0 \leqslant n < \mathrm{M}}\) 由置 \(e_n = f_{\iota(n)}\) （对 \(0 \leqslant n < \mathrm{M}\) ）在 E 中定义。它是 E 中的一个标准正交族。

在 M \textless{} Card(I) = +∞ 的情形，由 \(\{e_{0},\ldots,e_{M-1}\}\) 生成的空间 F 是有限维的，而它的正交 F° 是无穷维的；对 \((e_{n})_{n\geqslant M}\) 选取 F° 中的一个标准正交族。

对每个 \(x \in \mathrm{E}\) ，我们有

\[
u (x) = \sum_ {j \in J ^ {\prime}} \lambda_ {j} \langle f _ {j} \mid x \rangle f _ {j} = \sum_ {0 \leqslant n <   M} \lambda_ {n} (u) \langle e _ {n} \mid x \rangle e _ {n}
\]

因为 \(u(f_{j}) = 0\) （当 \(j \in J - J'\) ）。若 \(n \in I\) 满足 \(n \geqslant M\) ，我们有 \(\lambda_{n}(u) = 0\) ，由此得到命题的第一个公式。第二个公式由它得出。

命题 4. --- 设 u 是 E 的紧正自同态。设 \(f \in \mathcal{C}(\mathbf{R}_{+})\) 是连续递增函数，满足 \(f(0) = 0\) 。

自同态 \(f(u)\) 是紧且正的，且对每个 \(n \in \mathrm{I}_{\mathrm{E}}\) ，我们有 \(\lambda_n(f(u)) = f(\lambda_n(u))\) 。

由 III, p. 91, 命题 6, b) 与 I, p. 117, 命题 15, a)，自同态 \(f(u)\) 是紧且正的。\(f(u)\) 的谱是在 \(f\) 下 \(u\) 的谱的像（I, p. 111, 推论 2）。若 \(\lambda \in \mathrm{Sp}_{\mathrm{s}}(f(u))\) ，则 \(\lambda\) 的谱重数是 \(\mu \in \mathrm{Sp}(u)\) （满足 \(f(\mu) = \lambda\) ）的谱重数之和（III, p. 84, 推论 2）。由于 \(f\) 递增，序列 \((f(\lambda_n(u)))_{n \in \mathrm{I_E}}\) 是递减的。断言于是由序列 \((\lambda_n(f(u)))_{n \in \mathrm{I_E}}\) 的定义得出。

\subsection{特征值的变分刻画}

在本节中，我们假设 K = C。

设 \(u\) 是 E 的紧正自同态，\((\lambda_n(u))_{n\in \mathrm{I_E}}\) 是 \(u\) 的特征值递减序列。我们置 \(\mathrm{I} = \mathrm{I_E}\) 。

对 E 的每个闭子空间 F，我们记

\[
r _ {\mathrm{F}} (u) = \inf _ {x \in \mathrm{F} - \{0 \}} \frac {\langle x \mid u (x) \rangle}{\| x \| ^ {2}}, \quad \mathrm{R} _ {\mathrm{F}} (u) = \sup _ {x \in \mathrm{F} ^ {\circ} - \{0 \}} \frac {\langle x \mid u (x) \rangle}{\| x \| ^ {2}},
\]

其中下确界（相应地，上确界）是在 \([0, +\infty]\) 中取的。

对每个 \(n \in \mathbf{N}\) ，用 \(\mathcal{F}_n\) 表示 E 的向量子空间 \(\mathrm{F} \subset \mathrm{E}\) （维数为 \(n\) ）所成的集。我们称子空间 \(\mathrm{F} \in \mathcal{F}_n\) 适应于 \(u\) ，如果它有一个标准正交基 \((f_i)_{0 \leqslant i \leqslant n-1}\) ，使得 \(u(f_i) = \lambda_i(u)f_i\) （对 \(0 \leqslant i \leqslant n-1\) ）。

命题 5. --- 设 \(n \in I\) 。

\begin{enumerate}
\def\labelenumi{\alph{enumi})}
\item
  对每个子空间 \(\mathrm{F} \in \mathcal{F}_{n+1}\) （适应于 \(u\) 的），我们有 \(\lambda_n(u) = r_{\mathrm{F}}(u)\) ；
\item
  对每个子空间 \(F \in F_{n}\) （适应于 u 的），我们有 \(\lambda_{n}(u) = R_{\mathrm{F}}(u)\) 。
\end{enumerate}

设 \(F \in F_{n+1}\) 是适应于 u 的子空间，\((f_i)_{0 \leqslant i \leqslant n}\) 是 F 的一个标准正交基，使得 \(u(f_i) = \lambda_i(u)f_i\) （对所有 i）。对 F 中所有 x，我们有

\[
\langle x \mid u (x) \rangle = \sum_ {0 \leqslant i \leqslant n} \lambda_ {i} (u) | \langle f _ {i} \mid x \rangle | ^ {2} \geqslant \lambda_ {n} (u) \sum_ {0 \leqslant i \leqslant n} | \langle f _ {i} \mid x \rangle | ^ {2} = \lambda_ {n} (u) \| x \| ^ {2},
\]

其中等式当 \(x = f_{n}\) 时成立。这蕴含 \(r_{\mathrm{F}}(u) = \lambda_{n}(u)\) ，从而得到断言 (a)。

设 \(F \in F_{n}\) 是适应于 u 的子空间。闭子空间 \(F^{\circ}\) 非零（因为 \(n = \dim(F) < \dim(E)\) ）且在 u 下稳定；由 u 通过过渡到子空间得到的自同态 \(\widetilde{u}\) （\(F^{\circ}\) 的）是紧且正的。

有 \(\mathrm{Sp}(\widetilde{u})\subset \mathrm{Sp}(u)\) 。我们还有 \(\mathrm{Sp}(\widetilde{u})\subset [0,\lambda_n(u)]\) 。事实上，只需验证 \(\lambda \leqslant \lambda_{n}(u)\) （对所有 \(\lambda \in \mathrm{Sp}_{\mathrm{s}}(\widetilde{u})\) ）。数 \(\lambda\) 这时是 \(u\) 的特征值，故存在 \(j\in I\) 使 \(\lambda = \lambda_j(u)\) 。于是 \(u\) 的相对于 \(\lambda_{j}(u)\) 的特征空间不包含于 F，这蕴含 \(\lambda_{j}(u)\leqslant \lambda_{n}(u)\) 。

假设 \(\lambda_n(u) > 0\) 。\(u\) 的相对于 \(\lambda_n(u)\) 的特征空间于是不包含于 F，从而 \(\lambda_n(u)\) 属于 \(\widetilde{u}\) 的敏感谱。我们推出 \(\sup (\mathrm{Sp}(\widetilde{u})) = \lambda_n(u)\) 。若 \(\lambda_n(u) = 0\) ，由于 \(\widetilde{u}\) 的谱这时化为 \(\{0\}\) ，我们得到同样的结果。

此外，按定义我们有 \(\mathrm{R}_{\mathrm{F}}(u)=\mathrm{R}_{\{0\}}(\widetilde{u})\) ，最后按 I, p. 139, 命题 9, a)，有 \(\mathrm{R}_{\{0\}}(\widetilde{u})=\sup(\mathrm{Sp}(\widetilde{u}))\) 。断言 b) 得证。

命题 6. --- 对所有 \(n \in I\) ，我们有

\[
\lambda_ {n} (u) = \sup _ {\mathrm{F} \in \mathscr {F} _ {n + 1}} r _ {\mathrm{F}} (u) = \inf _ {\mathrm{F} \in \mathscr {F} _ {n}} \mathrm{R} _ {\mathrm{F}} (u).
\]

设 \((e_{n})_{n\in\mathrm{I}}\) 是具有 IV, p. 152 命题 3 的性质的标准正交族。对每个整数 \(n\) （满足 \(1\leqslant n<M+1\) ），设 \(F_{n}\in\mathcal{F}_{n}\) 是由 \((e_{0},\ldots,e_{n-1})\) 生成的 E 的 n 维子空间；按构造，空间 \(F_{n}\) 适应于 u。于是按命题 5，我们有

\[
\lambda_ {n} (u) = r _ {\mathrm{F} _ {n + 1}} (u) = \mathrm{R} _ {\mathrm{F} _ {n}} (u).\tag{2}
\]

设 \(n \in I\) 。设 \(F \in \mathcal{F}_{n+1}\) 。在 \(F\) 上的限制（到 \(F_n\) 上的正交投影的）不是单射，故存在 \(x \neq 0\) （\(F\) 中的，正交于 \(F_n\) 的）。

由于 \(x \in \mathrm{F}_{n}^{\circ}\) ，我们这时有（IV, p. 152, 命题 3）

\[
\begin{array}{l}\langle x\mid u(x)\rangle = \sum_{\substack{m\in \mathrm{I}\\ m\geqslant n}}\lambda_{m}(u)|\langle e_{m}\mid x\rangle |^{2}\\ \leqslant \lambda_{n}(u)\sum_{\substack{m\in \mathrm{I}\\ m\geqslant n}}|\langle e_{m}\mid x\rangle |^{2} = \lambda_{n}(u)\| x\|^{2}. \end{array}
\]

这证明 \(r_{\mathrm{F}}(u)\leqslant\lambda_{n}(u)\) ，从而特别地有不等式

\[
\sup _ {\mathrm{F} \in \mathscr {F} _ {n + 1}} r _ {\mathrm{F}} (u) \leqslant \lambda_ {n} (u).\tag{3}
\]

设 \(F \in F_{n}\) 。在 \(F_{n+1}\) 上的限制（到 F 上的正交投影的）不是单射，故存在向量 \(x \neq 0\) （\(F_{n+1}\) 中的，正交于 F 的）。由于 \(x \in F_{n+1}\) ，我们有（同处）

\[
\begin{array}{c} \langle x \mid u (x) \rangle = \sum_ {0 \leqslant m \leqslant n} \lambda_ {m} (u) | \langle e _ {m} \mid x \rangle | ^ {2} \\ \geqslant \lambda_ {n} (u) \sum_ {0 \leqslant m \leqslant n} | \langle e _ {m} \mid x \rangle | ^ {2} = \lambda_ {n} (u) \| x \| ^ {2} \end{array}
\]

因而 \(\mathrm{R_F}(u) \geqslant \lambda_n(u)\) 。特别地，由此得出

\[
\inf _ {\mathrm{F} \in \mathcal {F} _ {n}} \mathrm{R} _ {\mathrm{F}} (u) \geqslant \lambda_ {n} (u).\tag{4}
\]

鉴于公式 (2)、(3) 与 (4)，命题得证。

\subsection{特征值变分刻画的应用}

在本节中，我们考虑 C 上的希尔伯特空间。

命题 7. --- 设 u 与 v 是 E 的正紧自同态。

\begin{enumerate}
\def\labelenumi{\alph{enumi})}
\item
  有 \(|\lambda_n(u) - \lambda_n(v)|\leqslant \| u - v\|\) （对所有 \(n\in \mathbf{I}\) ）；
\item
  若 \(u \leqslant v\) ，则 \(\lambda_n(u) \leqslant \lambda_n(v)\) （对所有 \(n \in I\) ）。
\end{enumerate}

设 \(n \in I\) ，设 F 是 E 的 n 维向量子空间。对所有 \(x \in F\) ，我们有

\[
| \langle x | v (x) \rangle - \langle x | u (x) \rangle | \leqslant \| u - v \| \| x \| ^ {2},
\]

因此有不等式

\[
\mathrm{R} _ {\mathrm{F}} (v) - \| u - v \| \leqslant \mathrm{R} _ {\mathrm{F}} (u) \leqslant \mathrm{R} _ {\mathrm{F}} (v) + \| u - v \|.
\]

断言 \(a\) ) 由此按 IV, p. 154 命题 6 得出。

若 \(u \leqslant v\) ，我们有 \(\langle x | u(x) \rangle \leqslant \langle x | v(x) \rangle\) （对所有 \(x \in E\) ）。因此对每个 \(n \in I\) 与每个子空间 \(F\) （维数为 \(n\) 的），我们有 \(R_F(u) \leqslant R_F(v)\) ，从而 \(\lambda_n(u) \leqslant \lambda_n(v)\) （同处）。

命题 8. --- 设 u 是 E 的正紧自同态。设 H 是 E 的一个闭子空间，\(i_{\mathrm{H}}: \mathrm{H} \to \mathrm{E}\) 是典范单射。用 \(u_{\mathrm{H}}\) 表示 H 的自同态 \(i_{\mathrm{H}}^{*}ui_{\mathrm{H}}\) 。它是紧且正的。

\begin{enumerate}
\def\labelenumi{\alph{enumi})}
\item
  有 \(\mathrm{I_H}\subset \mathrm{I_E}\) 与 \(\lambda_n(u_{\mathrm{H}})\leqslant \lambda_n(u)\) （对所有 \(n\in \mathrm{I_H}\) ）；
\item
  若 H 在 E 中有有限余维数 \(k \in N\) ，则我们有 \(I_{H} + k \subset I_{E}\) 与 \(\lambda_{n+k}(u) \leqslant \lambda_{n}(u_{\mathrm{H}})\) （对所有 \(n \in I_{H}\) ）。
\end{enumerate}

自同态 \(u_{\mathrm{H}}\) 是紧的（III, p. 5, 命题 3）。它是正的，因为 \(\langle x \mid u_{\mathrm{H}}(x) \rangle = \langle i_{\mathrm{H}}(x) \mid u(i_{\mathrm{H}}(x)) \rangle \geqslant 0\) （对所有 \(x \in \mathrm{H}\) ）。

设 \(n \in I_{H} \subset I_{E}\) 。设 F 是 H 的一个维数为 \(n + 1\) 的适应于 \(u_{H}\) 的子空间。于是我们有 \(\lambda_{n}(u_{\mathrm{H}}) = r_{\mathrm{F}}(u_{\mathrm{H}})\) （IV, p. 153, 命题 5, a)），又由于 \(r_{\mathrm{F}}(u_{\mathrm{H}}) = r_{\mathrm{F}}(u) \leqslant \lambda_{n}(u)\) （IV, p. 154, 命题 6），我们得到断言 a)。

假设 H 在 E 中有余维数 \(k \in N\) 且 \(n \in I_{H}\) 。设 F 是 H 的一个维数为 n 的适应于 \(u_{H}\) 的子空间。它在 H 中的正交等于 \(H \cap F^{\circ}\) ，且它是 E 中子空间 \(F + H^{\circ}\) （维数为 \(n + k\) ）的正交。因此 \(n + k \in I_{E}\) ，且（IV, p. 153, 命题 5, b)）

\[
\lambda_{n}(u_{\mathrm{H}}) = \sup_{\substack{x\in \mathrm{H}\cap \mathrm{F}^{\circ}\\ x\neq 0}}\frac{\langle x\mid u_{\mathrm{H}}(x)\rangle}{\|x\|^{2}} = \mathrm{R}_{\mathrm{F} + \mathrm{H}^{\circ}}(u)
\]

从而 \(\lambda_{n}(u_{\mathrm{H}})\leqslant \lambda_{n + k}(u)\) （IV, p. 154, 命题 6）。

定义 3. --- 设 F 是希尔伯特空间，u 是从 E 到 F 中的一个紧线性映射。

对每个整数 \(n \in \mathrm{I}_{\mathrm{E}}\) ，置 \(\alpha_{n}(u) = \lambda_{n}(u^{*} \circ u)\) 。设 J 是由 \(n \in \mathrm{I}_{\mathrm{E}}\) （使 \(\alpha_{n}(u) > 0\) ）组成的集。族 \((\alpha_{n}(u))_{n \in \mathrm{J}}\) 称为 \(u\) 的按重数重复的奇异值序列。

我们称序列 \((\alpha_{n}(u))_{n\in\mathrm{I}_{\mathrm{E}}}\) 为 u 的增广奇异值序列。

序列 \((\alpha_{n}(u))_{n\in\mathrm{I}_{\mathrm{E}}}\) 是完全定义的，因为 E 的自同态 \(u^{*}\circ u\) 是紧的（III, p. 5, 命题 3）；由于 \(u^{*}\circ u\) 是正的，它是一个正实数的递减族。

命题 9. --- 设 F 是希尔伯特空间，u 是从 E 到 F 中的一个紧线性映射。

\begin{enumerate}
\def\labelenumi{\alph{enumi})}
\tightlist
\item
  对 \(n \in \mathrm{I}_{\mathrm{E}}\) ，我们有
\end{enumerate}

\[
\alpha_ {n} (u) = \sup _ {\mathrm{F} \in \mathscr {F} _ {n + 1}} \inf _ {x \in \mathrm{F} - \{0 \}} \frac {\| u (x) \|}{\| x \|} = \inf _ {\mathrm{F} \in \mathscr {F} _ {n}} \sup _ {x \in \mathrm{F} ^ {\circ} - \{0 \}} \frac {\| u (x) \|}{\| x \|}.
\]

\begin{enumerate}
\def\labelenumi{\alph{enumi})}
\setcounter{enumi}{1}
\tightlist
\item
  设 J 是由 \(n \in I_{E}\) （使 \(\alpha_{n}(u) \neq 0\) ）组成的集。存在 E 中的标准正交族 \((e_{n})_{n \in \mathrm{J}}\) 与 F 中的标准正交族 \((f_{n})_{n \in \mathrm{J}}\) ，使得对每个 \(x \in E\) ，我们有
\end{enumerate}

\[
u (x) = \sum_ {n \in J} \alpha_ {n} (u) \langle e _ {n} | x \rangle f _ {n}.
\]

由于 \(\langle x \mid u^{*}u(x)\rangle = \|u(x)\|^{2}\) （对所有 \(x \in E\) ），诸定义与 IV, p. 154, 命题 6, a) 蕴含第一个断言中的等式。设 \((e_{n})_{n\in\mathrm{I}_{\mathrm{E}}}\) 是 E 的一个标准正交族，满足把 IV, p. 152 命题 3 应用于 E 的紧正自同态 \(u^{*} \circ u\) 所得的结论。置 \(f_{n} = \alpha_{n}^{-1}u(e_{n})\) （对 \(n \in J\) ）。按 IV, p. 150 推论 2 的证明中的推理，我们得到断言 c)。

推论. --- 设 F 是希尔伯特空间，u 是从 E 到 F 中的一个紧线性映射。设 \(w \in \mathcal{L}(\mathrm{E})\) 与 \(v \in \mathcal{L}(\mathrm{F})\) 。对每个 \(n \in I_{E}\) ，我们有 \(\alpha_{n}(w \circ u \circ v) \leqslant \|v\|\) \(\|w\|\) \(\alpha_{n}(u)\) 。

这是前一命题断言 \(a\) ) 的推论。

注记. --- 1) 序列 \((\alpha_{n}(u))_{n\in\mathrm{I}_{\mathrm{E}}}\) 是递减的，故集 J 或等于 \(\mathrm{I}_{\mathrm{E}}\) ，或等于线段 \(\{0,\ldots,m\}\) （\(\mathrm{I}_{\mathrm{E}}\) 中的）（其中 \(m\in\mathbf{N}\) ）。后一情形成立当且仅当 \(u\) 有限秩。

\begin{enumerate}
\def\labelenumi{\arabic{enumi})}
\setcounter{enumi}{1}
\item
  由于 \(\| u(x)\| = \| |u|(x)\|\) （对所有 \(x\in \mathrm{E}\) ；I, p. 139, 命题 10），我们有 \(\alpha_{n}(u) = \alpha_{n}(|u|)\) （对所有 \(n\in \mathrm{I}_{\mathrm{E}}\) ）。
\item
  若 \(u\) 是正的，则 \(\alpha_{n}(u)=\lambda_{n}(u)\) （对所有 \(n\in\mathrm{I}_{\mathrm{E}}\) ；事实上，按 IV, p. 153 命题 4，我们这时有 \(\alpha_{n}(u)^{2}=\lambda_{n}(u^{*}u)=\lambda_{n}(u^{2})=\lambda_{n}(u)^{2}\) ）。
\item
  有可能有 \(\alpha_{n}(v\circ u\circ w)=0\) ，即使 \(\alpha_{n}(u)\) 非零；在这种情形，\(\alpha_{n}(v\circ u\circ w)\) 不是 \(v\circ u\circ w\) 的奇异值。
\end{enumerate}